\documentclass[11pt]{article}

\usepackage[a4paper,margin=1in]{geometry}
\usepackage{amsmath,amssymb,amsthm,mathtools}
\usepackage{booktabs,array,longtable}
\usepackage{enumitem}
\usepackage{microtype}
\usepackage{xcolor}
\usepackage{authblk}
\usepackage{hyperref}
\usepackage[nameinlink,capitalise,noabbrev]{cleveref}
\renewcommand{\Affilfont}{\small}
\renewcommand{\Authfont}{\normalsize}


\hypersetup{
  colorlinks=true,
  linkcolor=blue!55!black,
  citecolor=blue!55!black,
  urlcolor=blue!55!black
}

\newtheorem{theorem}{Theorem}[section]
\newtheorem{proposition}[theorem]{Proposition}
\newtheorem{lemma}[theorem]{Lemma}
\newtheorem{corollary}[theorem]{Corollary}
\newtheorem{definition}[theorem]{Definition}
\newtheorem{remark}[theorem]{Remark}

% The environments above share the theorem counter, so cleveref would name
% them all "Theorem"; these aliases restore the correct names.
\usepackage{etoolbox}
\AtBeginEnvironment{proposition}{\crefalias{theorem}{proposition}}
\AtBeginEnvironment{lemma}{\crefalias{theorem}{lemma}}
\AtBeginEnvironment{corollary}{\crefalias{theorem}{corollary}}
\AtBeginEnvironment{definition}{\crefalias{theorem}{definition}}
\AtBeginEnvironment{remark}{\crefalias{theorem}{remark}}

\numberwithin{equation}{section}

\newcommand{\one}{\mathbf 1}
\newcommand{\R}{\mathbb R}
\newcommand{\ip}[2]{\left\langle #1,#2\right\rangle}
\newcommand{\norm}[1]{\left\lVert #1\right\rVert}
\newcommand{\Tr}{\operatorname{Tr}}
\newcommand{\ord}{\operatorname{ord}}
\newcommand{\pos}[1]{\left(#1\right)_{+}}
\newcommand{\dd}{\,\mathrm d}
\newcommand{\E}{\mathbb E}
\newcommand{\BetaLaw}{\operatorname{Beta}}
\newcommand{\supp}{\operatorname{supp}}

\newcommand*{\tk}[1]{\textcolor{blue}{\textbf{[TK: #1]}}}
\newcommand*{\jl}[1]{\textcolor{red!70!black}{\textbf{[JL: #1]}}}
\newcommand*{\jb}[1]{\textcolor{green!50!black}{\textbf{[JB: #1]}}}
\newcommand*{\si}[1]{\textcolor{orange!80!black}{\textbf{[SI: #1]}}}
\newcommand*{\jk}[1]{\textcolor{purple}{\textbf{[JK: #1]}}}
\newcommand*{\hy}[1]{\textcolor{teal}{\textbf{[HY: #1]}}}

\newcommand{\dens}[2]{t(#1,#2)}
\newcommand{\tC}[2]{t(C_{#1},#2)}
\newcommand{\tP}[2]{t(P_{#1},#2)}
\newcommand{\edens}[1]{t(K_{2},#1)}

\newcommand{\goodman}[2]{#1^{#2}-#1(1-#1)^{#2-1}}
\newcommand{\goal}{\goodman{p}{m}}

\newcommand{\Tker}[1]{T_{#1}}
\newcommand{\dW}{d_W}

\newcommand{\coeff}[1]{[z^{#1}]}
\newcommand{\nld}[1]{-\log\det(I-z#1)}

\newcommand{\resm}{(I-zA)^{-1}}
\newcommand{\resp}{(I+zA)^{-1}}
\newcommand{\gresm}{\ip g{\resm g}}
\newcommand{\gresp}{\ip g{\resp g}}

\newcommand{\specmu}{\sum_i\norm{P_ig}^2\,\delta_{\lambda_i}}
\newcommand{\dmu}{\dd\mu(\lambda)}
\newcommand{\intmu}[1]{\int #1\dmu}
\newcommand{\gA}[1]{\ip g{A^{#1}g}}

\newcommand{\intsq}{\int_{[0,1]^2}}
\newcommand{\intcube}[1]{\int_{[0,1]^{#1}}}
\newcommand{\ddxy}{\dd x\dd y}

\newcommand{\rep}[2]{#1^{\times #2}}

\newcommand{\blockmat}[4]{\begin{pmatrix}#1&#2\\#3&#4\end{pmatrix}}
\newcommand{\sblockmat}[4]{\left(\begin{smallmatrix}#1&#2\\#3&#4\end{smallmatrix}\right)}

\title{Goodman-type lower bounds for odd cycles I: The dense regime}
\author[1,2]{Taeyoung~Kim}
\author[3]{Jineon~Baek}
\author[4]{Seonghyuk~Im}
\author[2]{Joonkyung~Lee}
\author[1,2]{Hongseok~Yang}
\affil[1]{School of Computing, KAIST, Korea}
\affil[2]{School of Computational Sciences, Korea Institute for Advanced Study (KIAS), Korea}
\affil[3]{June E Huh Center for Mathematical Challenges, Korea Institute for Advanced Study (KIAS), Korea}
\affil[4]{Center for AI and Natural Science, Korea Institute for Advanced Study (KIAS), Korea}
\affil[ ]{\newline\texttt{taeyoungkim21@kaist.ac.kr}, \texttt{{jineon.kias}@gmail.com}, \texttt{seonghyuk@kias.re.kr}, \texttt{joonkyunglee@kias.re.kr}, \texttt{hongseokyang@kias.re.kr}}
\date{\today}

\begin{document}
\maketitle

\begin{abstract}
  A graph is bipartite exactly when it contains no odd cycle.
  This characterization makes odd-cycle counts a natural quantitative measure of departure from bipartiteness and leads to the problem of determining how many odd cycles are forced by a prescribed edge density.
  Goodman's triangle inequality provides a well-known lower bound in the triangle case.
  We establish its analogue for every odd cycle in the dense regime: in the language of graph limits, for every odd integer $m\ge3$, every graphon $W$ with $\edens W=p\ge2/3$ satisfies
  \[
    \tC mW\ge p^m-p(1-p)^{m-1}.
  \]
  Here $\tC mW$ denotes the $C_m$-density of $W$.
  At every density $p=1-1/k$ with an integer $k\ge2$, the balanced complete $k$-partite graphon attains equality.
  For $k\ge3$, we also obtain uniqueness up to graphon equivalence and qualitative stability.
  Our theorem extends the $C_5$ bound of Bennett, Dudek, Lidick\'{y} and Pikhurko~\cite{bennett2020minimizing} to every odd cycle.
  This is the first of two papers establishing the bound for all edge densities.
  A companion paper treats $1/2<p<2/3$; for $p\le1/2$, the inequality is immediate because its right-hand side is nonpositive.
  At every edge density, we also obtain a density-weighted strengthening of Kim and Lee's extended-commonality inequality~\cite{kim2024extended} for odd cycles.
  Our proof combines spectral cancellation, cycle and path excursion expansions, Dirichlet/Beta averaging, and elementary spectral moment estimates.
  The density-weighted comparison has a positive expectation representation at every edge density; a separate comparison of complementary even-path corrections yields the dense-regime bound.
  Together, these tools give a single argument that is uniform in $m$, with an explicit degree-variance remainder.
  The one-colour lower-bound theorem has also been formally verified in Lean~4 by a separate argument; the code is available at \url{ADDLINKLATER}.
\end{abstract}

\section{Introduction}\label{sec:intro}

How does a graph's global edge density force specific local configurations?
Extremal graph theory first addresses the threshold form of this question: how dense can a graph be while avoiding a prescribed subgraph $H$?
The Erd\H os--Stone--Simonovits theorem~\cite{erdos1966limit,erdos1946structure} answers it asymptotically: if $H$ has chromatic number $\chi(H)=r+1$, then every $H$-free graph has edge density at most $1-1/r+o(1)$.
The supersaturation theorem of Erd\H os and Simonovits~\cite{erdos1983supersaturated} further shows that exceeding this threshold by a fixed amount forces a positive density of copies of $H$.
These results identify when copies of $H$ become unavoidable, but leave open the quantitative problem studied here: among graphs of edge density $p$, how small can the density of $H$ be?

Odd cycles are a natural family for this problem because a graph is bipartite exactly when it contains no odd cycle.
Their densities therefore provide a quantitative measure of departure from bipartiteness, with a threshold at $p=1/2$.
Bipartite graphs can have any asymptotic edge density $p\le1/2$ while avoiding all odd cycles.
Once $p>1/2$, the theorems above force a positive density of every fixed odd cycle.
We seek an explicit lower bound for these densities throughout this range.

For complete subgraphs, the corresponding question is the classical Erd\H os--Rademacher, or clique-density, problem~\cite{erdos1962number}.
Goodman~\cite{goodman1959sets} found a simple polynomial lower bound for triangles that is attained by the balanced complete $k$-partite graph at each edge density $p=1-1/k$, $k\ge2$, but is not generally sharp between these densities.
The exact minimum follows a more intricate, piecewise-defined curve.
Razborov~\cite{razborov2008minimal} determined this curve for triangles using flag algebras, Nikiforov~\cite{nikiforov2010number} treated $K_4$, and Reiher~\cite{reiher2016clique} proved the clique density theorem for every complete graph, establishing the asymptotic minimum conjectured by Lov\'asz and Simonovits~\cite{lovasz1983number}.
The asymptotic structure of extremal and almost extremal graphs was characterized by Pikhurko and Razborov~\cite{pikhurko2017asymptotic} for triangles and by Kim, Liu, Pikhurko and Sharifzadeh~\cite{kim2020asymptotic} for general cliques.

For odd cycles, our aim is an analogue of Goodman's polynomial bound: a simple formula that works uniformly for every odd cycle length, with no case distinctions in $p$.
We seek equality at each density $p=1-1/k$, where the balanced complete $k$-partite graph supplies the candidate extremal value, while allowing the bound to be nonoptimal between these densities.

We formulate the problem using graphons, which encode limits of dense graphs and express asymptotic edge and subgraph densities as analytic quantities~\cite{lovasz2012large,lovasz2006limits}.
A \emph{graphon} is a symmetric measurable function $W\colon[0,1]^2\to[0,1]$.
For a finite graph $H$ with vertex set $\{1,\ldots,v\}$ and edge set $E(H)$, its \emph{homomorphism density} in $W$ is
\[
  \dens HW=\intcube{v}\prod_{ij\in E(H)}W(x_i,x_j)\dd x_1\cdots\dd x_v,
\]
and the \emph{edge density} of $W$ is $\edens W=\intsq W(x,y)\ddxy=p$.
Thus, for an odd integer $m\ge3$ and a fixed edge density $p\in[0,1]$, we seek a simple lower bound for
\[
  \inf\bigl\{\tC mW: W\text{ is a graphon with }\edens W=p\bigr\},
\]
where $C_m$ denotes the cycle of length $m$.
The spectral characterization of cycle densities suggests a candidate bound.
Associate with a graphon $W$ the integral operator $\Tker W$ on $L^2[0,1]$ defined by
\[
  (\Tker{W}f)(x)=\int_0^1 W(x,y)f(y)\dd y.
\]
Because $W$ is symmetric, $\Tker W$ is self-adjoint.
If $(\lambda_i)_i$ are its eigenvalues, then the trace formula~\cite{lovasz2012large} gives
\[
  \tC mW=\Tr(\Tker W^m)=\sum_i\lambda_i^m.
\]
Thus, cycle densities are power sums of the eigenvalues of $\Tker W$.

For balanced complete multipartite graphons, these eigenvalues can be computed explicitly.
For an integer $k\ge2$, let $W_k$ denote the balanced complete $k$-partite graphon defined by
\[
  W_k(x,y)=
  \begin{cases}
    1&\text{if }\lfloor kx\rfloor\ne\lfloor ky\rfloor,\\
    0&\text{otherwise}.
\end{cases}
\]
Its edge density is $p=1-1/k$.
Write $I_1,\ldots,I_k$ for the $k$ equal parts.
If $f$ takes the constant value $c_i$ on each part $I_i$, then
\[
  (\Tker{W_k}f)(x)=\sum_{j\ne i}\int_{I_j}f(y)\dd y
  =\frac1k\sum_{j=1}^k c_j-\frac{c_i}{k},\qquad x\in I_i.
\]
Thus, the constant function has eigenvalue $(k-1)/k=p$, while the partwise constant functions with $\sum_i c_i=0$ form a $(k-1)$-dimensional eigenspace with eigenvalue $-1/k$.
Every function decomposes into its partwise averages and a remainder whose integral on each part is zero.
The operator maps this remainder to zero, so all remaining eigenvalues are zero.
The trace formula therefore gives, for odd $m$,
\begin{equation}\label{eq:intro-kpartite-value}
  \tC m{W_k}=p^m+(k-1)\Bigl(-\frac1k\Bigr)^m=\goal.
\end{equation}
At each density $p=1-1/k$, this construction gives an upper bound on the minimum $C_m$-density.
Motivated by this expression and Goodman's triangle inequality, we ask whether $\goal$ is also a lower bound for every graphon.
Such a bound would establish the optimality of $W_k$ at these densities, while allowing the bound to be nonoptimal between them.

In the dense range $p\ge2/3$ considered here, this polynomial bound was previously known only for $C_3$ and $C_5$.
For $m=3$, it is exactly Goodman's inequality: every graphon $W$ with $\edens W=p$ satisfies
\begin{equation}\label{eq:intro-goodman}
  \tC3W\ge2p^2-p=\goodman p3.
\end{equation}
For $m=5$, Bennett, Dudek, Lidick\'y and Pikhurko~\cite{bennett2020minimizing} proved that, for every integer $k\ge3$, $W_k$ minimizes the $C_5$-density among graphons $W$ with $\edens W=1-1/k$.
They establish the lower bound by a flag algebra computation and also prove a stability theorem describing graphs whose $C_5$-density is close to the minimum.
Although their theorem is stated only for integer $k\ge3$, a direct inspection of their flag algebra certificate shows that it remains valid for every real $k\ge3$.
Substituting $k=1/(1-p)$ yields $\tC5W\ge\goodman p5$ for $2/3\le p<1$; the case $p=1$ is immediate.

Our main result extends this bound to every odd cycle in the dense range.

\begin{theorem}[Odd-cycle Goodman bound for $p\ge2/3$]\label{thm:main}
  Let $m\ge3$ be odd, and let $W$ be a graphon with $\edens W=p\ge2/3$.
  Then,
  \begin{equation}\label{eq:intro-main}
    \tC mW\ge\goal.
  \end{equation}
\end{theorem}

This is the first of two papers establishing the bound at every edge density.
A companion paper treats $1/2<p<2/3$, and the case $p\le1/2$ is immediate because $p^m-p(1-p)^{m-1}$ is nonpositive.
At the edge densities $p=1-1/k$ with $k\ge3$, \Cref{thm:main} yields the exact extremal value.

\begin{corollary}\label{cor:kpartite}
  For every integer $k\ge2$ and every odd $m\ge3$, the balanced complete $k$-partite graphon $W_k$ minimizes the $C_m$-density among graphons $W$ with $\edens W=1-1/k$.
  For $k\ge3$, it is the unique minimizer up to graphon equivalence.
\end{corollary}

For $k\ge3$, \eqref{eq:intro-kpartite-value} shows that $W_k$ attains the lower bound in \Cref{thm:main}.
For $k=2$, the conclusion follows because $W_2$ is bipartite, so $\tC m{W_2}=0$, while every graphon has nonnegative $C_m$-density.
The uniqueness assertion follows from the equality analysis in \Cref{sec:equality-stability}, where graphon equivalence is defined by zero cut distance.
It also gives qualitative stability: at each fixed density $1-1/k$ with $k\ge3$, graphons whose $C_m$-densities approach the minimum must approach $W_k$ in cut distance (\Cref{cor:stability}).

\paragraph{A stronger two-colour theorem.}
Our proof of \Cref{thm:main} passes through a stronger inequality involving $W$ and its complement $1-W$.
This inequality strengthens a density-sensitive form of \emph{commonality} for odd cycles.
A graph $H$ with $e(H)$ edges is common if
\[
  \dens HW+\dens H{1-W}\ge2^{1-e(H)}
\]
for every graphon $W$.
Here $W$ and $1-W$ represent the two colours; the inequality says that the uniform random colouring asymptotically minimizes the number of monochromatic copies of $H$.
Goodman's theorem shows that the triangle is common.
Thomason~\cite{thomason1989disproof} showed that $K_4$ is not common, disproving conjectures of Erd\H os~\cite{erdos1962number} and of Burr and Rosta~\cite{burr1980ramsey}; Jagger, \v{S}\v{t}ov\'{\i}\v{c}ek and Thomason~\cite{jagger1996multiplicities} later proved that no graph containing $K_4$ is common.
Sidorenko~\cite{sidorenko1989cycles} proved that every cycle is common.

Kim and Lee~\cite{kim2024extended} established the following density-sensitive strengthening, called \emph{extended commonality}, for every path or cycle $H$ and every graphon $W$ with $\edens W=p$:
\begin{equation}\label{eq:intro-extended-commonality}
  \dens HW+\dens H{1-W}\ge p^{e(H)}+(1-p)^{e(H)}.
\end{equation}
This answers a question of Sidorenko~\cite{sidorenko1989cycles}.
It also proves the conjecture that every odd cycle is strongly common, posed by Behague, Morrison and Noel in the first arXiv version of their paper~\cite[Conjecture~9.7]{behague2022commonv1}.
Since $p^{e(H)}+(1-p)^{e(H)}\ge2^{1-e(H)}$, extended commonality implies commonality.

For odd cycles, our second main theorem strengthens
\eqref{eq:intro-extended-commonality} at every edge density by a
symmetric, density-weighted path correction. Throughout, $P_k$ denotes
the path with $k$ edges.

\begin{theorem}[Density-weighted extended commonality]
\label{thm:weighted-commonality}
  Let $m\ge3$ be odd, and let $W$ be a graphon with
  $\edens W=p$ and $q=1-p$. Then
  \begin{equation}\label{eq:intro-weighted-commonality}
    \begin{aligned}
      \tC mW+\tC m{1-W}\ge{}&p^m+q^m\\
      &+\frac{m}{m-1}p\bigl(\tP{m-1}W-p^{m-1}\bigr)\\
      &+\frac{m}{m-1}q\bigl(\tP{m-1}{1-W}-q^{m-1}\bigr).
    \end{aligned}
  \end{equation}
  Equality holds if and only if $W$ has constant degree $p$, that is,
  $\int_0^1W(x,y)\dd y=p$ for almost every $x$.
\end{theorem}

The Blakley--Roy inequality~\cite{blakley1965holder} (independently
proved by Mulholland and Smith~\cite{mulholland1959inequality}) gives
$\tP{m-1}W\ge p^{m-1}$ and
$\tP{m-1}{1-W}\ge q^{m-1}$. Thus both path corrections are
nonnegative, and \Cref{thm:weighted-commonality} implies the
odd-cycle case of \eqref{eq:intro-extended-commonality}. Its right-hand
side is strictly larger than $p^m+q^m$ for every nonregular graphon;
we verify this using even path moments in \Cref{sec:dense-paths}.
The weights $p$ and $q$ vary continuously with the edge density, and
\eqref{eq:intro-weighted-commonality} is unchanged by interchanging
$W$ and $1-W$.

The proof also gives a form that retains an explicit degree-variance
term:
\begin{equation}\label{eq:intro-weighted-variance}
  \begin{aligned}
    \tC mW+\tC m{1-W}\ge{}&
      p\,\tP{m-1}W+q\,\tP{m-1}{1-W}\\
      &+(m-1)\bigl(\tP2W-p^2\bigr)^{(m-1)/2}.
  \end{aligned}
\end{equation}
This inequality also holds at every edge density. The two formulations
are obtained from the same positive expectation identity by two
choices of its coefficient, as shown in
\Cref{prop:weighted-variance}. The coefficient-one formulation
\eqref{eq:intro-weighted-variance} retains the variance remainder
needed for the dense-regime equality analysis.

To deduce the one-colour bound, we compare the two even-path
corrections. In \Cref{lem:dense-path-comparison} we prove
\begin{equation}\label{eq:intro-path-comparison}
  \tP kW-p^k\ge\tP k{1-W}-q^k
  \qquad\text{for every even $k\ge2$, if $p\ge2/3$.}
\end{equation}
Combining this comparison with
\eqref{eq:intro-weighted-variance} yields the following dense-regime
consequence.

\begin{corollary}[Complementary-path correction in the dense regime]
\label{thm:commonality}
  Let $m\ge3$ be odd, and let $W$ be a graphon with
  $\edens W=p\ge2/3$. Set $q=1-p$. Then
  \begin{equation}\label{eq:intro-commonality}
    \tC mW+\tC m{1-W}\ge p^m+q^m+
       \bigl[\tP{m-1}{1-W}-q^{m-1}\bigr].
  \end{equation}
\end{corollary}

Indeed, applying \eqref{eq:intro-path-comparison} with $k=m-1$
and using $p+q=1$ gives
\begin{align*}
  &p\bigl(\tP{m-1}W-p^{m-1}\bigr)
    +q\bigl(\tP{m-1}{1-W}-q^{m-1}\bigr)\\
  &\qquad\ge
    (p+q)\bigl(\tP{m-1}{1-W}-q^{m-1}\bigr)
    =\tP{m-1}{1-W}-q^{m-1}.
\end{align*}
After the baseline terms are restored, this proves
\eqref{eq:intro-commonality}, and retaining the variance term gives
\Cref{prop:regularity-gap}. Equality in
\eqref{eq:intro-commonality} holds exactly when $W$ has constant
degree $p$. This includes the constant graphon $W\equiv p$ and,
when $p=1-1/k$ with $k\ge3$, $W_k$. Exchanging $W$ and $1-W$
gives the corresponding complementary-path statement for $p\le1/3$.

The dense-regime statement immediately implies \Cref{thm:main}.
Since $0\le1-W\le1$, deleting one edge from the $m$-cycle gives
$\tC m{1-W}\le\tP{m-1}{1-W}$. Rearranging
\eqref{eq:intro-commonality} and using
$q^m-q^{m-1}=-pq^{m-1}$, we obtain
\[
  \tC mW-\bigl(\goal\bigr)
  \ge\tP{m-1}{1-W}-\tC m{1-W}\ge0.
\]

\paragraph{Proof strategy: spectral cancellation, averaging, and path comparison.}
The trace formula expresses the cycle densities as
$\tC mW=\Tr(\Tker W^m)$ and
$\tC m{1-W}=\Tr(\Tker{1-W}^m)$. We separate constant functions
from zero-average functions. More precisely, for $f\in L^2[0,1]$,
write $f=c\one+f_0$, where $c=\int_0^1f$ and $\int_0^1f_0=0$,
and define the orthogonal projections
\[
  \Pi_{\mathrm c}f=c\one,
  \qquad \Pi_0f=f_0.
\]
For either $T=\Tker W$ or $T=\Tker{1-W}$, insertion of
$\Pi_{\mathrm c}+\Pi_0=I$ in the trace separates three types of terms:
those using only constants, those using only zero-average functions,
and those visiting both subspaces. The first contributions sum to
$p^m+q^m$. The second contributions cancel when $m$ is odd, since
$\Tker W+\Tker{1-W}$ is the projection onto constants.

The remaining terms can be organized by the number $r$ of passages
from the constant subspace into its orthogonal complement and back.
We call these passages \emph{excursions}. For a step graphon, the
block decomposition of the two operators defines a finite nonnegative
spectral measure $\mu$ with
\[
  \mu(\R)=\tP2W-p^2.
\]
The cycle and path excursion expansions express the difference
between the two sides of the density-weighted inequality as a sum of
integrals of complete homogeneous symmetric polynomials against
$\mu^{\otimes r}$. The same-density step approximation then reduces
the general graphon case to these finite-matrix identities.

The choice of density weights is important. Multiplication of a path
term by its edge density uses the elementary recurrence
\[
  a\,h_{n-1}(a^{\times(r+1)},\boldsymbol b)
  =h_n(a^{\times(r+1)},\boldsymbol b)
   -h_n(a^{\times r},\boldsymbol b).
\]
It converts the cycle--path comparison into a comparison of two
Dirichlet averages with adjacent parameters. If
$X\sim\BetaLaw(r,r)$ and $L$ is an independent uniformly weighted
convex combination of the spectral variables, the two linear forms
\[
  pX-(1-X)L,
  \qquad qX+(1-X)L
\]
sum to $X>0$. Their $(m-2r)$th powers therefore have positive sum,
because $m-2r$ is odd. For the coefficient $m/(m-1)$ in
\Cref{thm:weighted-commonality}, the remaining multiplier is a
positive constant times
\[
  r-1+m(1-X)>0.
\]
Thus each weighted excursion integrand is positive for every real
tuple of spectral variables. The coefficient-one version gives
\eqref{eq:intro-weighted-variance} by retaining its top excursion
term. These arguments require no restriction on $p$.

The density assumption enters only in the comparison
\eqref{eq:intro-path-comparison}. For paths with two or four edges,
this follows directly from the complementary block matrices and the
bound $\|A\|\le1/2$ for the zero-average compression of
$\Tker{1-W}$. For longer even paths, Jensen's inequality bounds the
path correction for $W$ from below by its degree variance, while the
operator-norm estimate
$\|\Tker{1-W}\|^2\le q$ bounds the complementary correction from
above by the same quantity. An elementary finite-sum estimate makes
the two bounds comparable when $p\ge2/3$. This proves
\Cref{thm:commonality}, after which edge deletion gives
\Cref{thm:main}.

\paragraph{Organization.}
\Cref{sec:matrix-reduction} sets up the finite-matrix representation
and identifies the spectral quantities used in the proof.
\Cref{sec:excursions} develops the cycle and path expansions, including
the density-weighted comparison. \Cref{sec:short-cycles} retains the
direct polynomial proof of the complementary-path bound for
$C_3,C_5,C_7$ in the larger range $p\ge1/2$.
\Cref{sec:excursion-positivity} proves the all-density weighted
inequalities by Dirichlet/Beta averaging, establishes the dense
complementary-path comparison, and completes the main proofs.
\Cref{sec:equality-stability} identifies the equality cases and
deduces qualitative stability. We close with further questions in
\Cref{sec:conclusion}.

\section{Matrix reduction and spectral measure}
\label{sec:matrix-reduction}

Let $W\colon[0,1]^2\to[0,1]$ be a graphon with edge density
$\edens W=p\in[0,1]$. Write
\[
  U=1-W,
  \qquad
  q=1-p.
\]
Thus, $U$ is the complementary graphon and has edge density $q$.
The assumption $p\ge2/3$ will be needed only for the complementary even-path comparison in \Cref{sec:dense-paths}. The matrix identities, excursion expansions, and density-weighted inequalities hold at every edge density.

To prove \Cref{thm:weighted-commonality,thm:commonality,thm:main} by finite matrix calculations, we first restrict attention to \emph{step graphons}: graphons that are constant on each product of two parts of a finite measurable partition of $[0,1]$.
We will use partitions into intervals of equal length, so these products form a grid of equally sized rectangles.
Each such graphon is specified by the finite matrix of its values, and its cycle and path densities can be expressed through matrix products.
The following lemma justifies this restriction. It applies to all the cycle and path inequalities below, including their degree-variance terms.

\begin{lemma}[Step-graphon reduction]
\label{lem:step-reduction}
  Every graphon $W$ admits a sequence of graphons, each constant on
  the rectangles of a partition into equal-length intervals, with
  exactly the same edge density as $W$, converging to $W$ in
  $L^1([0,1]^2)$. Their homomorphism densities, and those of their
  complements, converge for every finite graph. Consequently, it
  suffices to prove the inequalities in
  \Cref{thm:weighted-commonality,thm:commonality,thm:main} for such
  step graphons. The same conclusion applies to inequalities
  involving degree variance $\tP2W-p^2$.
\end{lemma}
Starting from an arbitrary graphon $W$, divide $[0,1]$ into $N$ equal intervals and replace $W$ on each rectangle $I_a\times I_b$ by its average value there.
This produces a step graphon with exactly the same edge density as $W$.
Starting with two equal intervals and repeatedly halving each partition interval gives increasingly fine grids, with $N=2,4,8,\ldots$ parts.
Averaging $W$ on each of these grids produces a sequence of step graphons whose cycle and path densities, and those of their complements, converge to the corresponding densities of $W$ and $1-W$.
Thus, each of the stated inequalities passes from this sequence to $W$ by taking limits.
The proof is given in \Cref{app:step-reduction}.

We may therefore fix such a step graphon with partition intervals
$I_1,\ldots,I_N$.
A \emph{step kernel} on this partition is a symmetric function $K\colon[0,1]^2\to\R$ that is constant on each rectangle $I_a\times I_b$.
It is a step graphon when its values lie in $[0,1]$; in particular, $W$, $U$, and the constant-one graphon are step kernels on this partition.
Write $K_{ab}$ for the value of $K$ on $I_a\times I_b$, and define the $N\times N$ matrix
$M_K$ and the vector $e\in\R^N$ by
\begin{equation*}
  M_K=\left(\frac{K_{ab}}N\right)_{a,b=1}^N,
  \qquad
  e=\frac1{\sqrt N}(1,\ldots,1)^{\mathsf T}.
\end{equation*}
Then, $M_K$ is symmetric and $\norm e=1$. For a step graphon $K$, this normalization expresses the relevant graph densities through matrix powers:
\begin{align}
  \Tr(M_K^j)
  &=\frac1{N^j}\sum_{a_1,\ldots,a_j=1}^N
    K_{a_1a_2}\cdots K_{a_ja_1}
   =\dens{C_j}K \qquad (j\ge3),\notag\\
  e^{\mathsf T}M_K^ke
  &=\frac1{N^{k+1}}\sum_{a_0,\ldots,a_k=1}^N
    K_{a_0a_1}\cdots K_{a_{k-1}a_k}
   =\dens{P_k}K \qquad (k\ge1),\notag\\
  e^{\mathsf T}M_Ke
  &=\frac1{N^2}\sum_{a,b=1}^NK_{ab}
   =\edens K. \label{eq:edge-matrix}
\end{align}

Let $\mathcal E_0=e^\perp$, and let $I$ be the $N\times N$ identity
matrix. The matrices
\[
  P=ee^{\mathsf T},
  \qquad
  Q=I-P
\]
are the orthogonal projections of $\R^N$ onto $\R e$ and $\mathcal E_0$,
respectively. They are the matrix counterparts of the projections
$\Pi_{\mathrm c}$ and $\Pi_0$ used in the introduction.
Relative to the decomposition $\R e\oplus\mathcal E_0$, these projections have the block forms
\begin{equation}\label{eq:block-PQ}
  P=\begin{pmatrix}1&0\\0&0\end{pmatrix},
  \qquad
  Q=\begin{pmatrix}0&0\\0&I_{\mathcal E_0}\end{pmatrix},
\end{equation}
where $I_{\mathcal E_0}$ is the identity on $\mathcal E_0$.

Define the vector $g\in\R^N$ and the symmetric linear map
$A\colon\mathcal E_0\to\mathcal E_0$ by
\begin{equation*}
  g=M_Ue-qe,
  \qquad
  A=QM_UQ\big|_{\mathcal E_0}.
\end{equation*}
The edge identity \eqref{eq:edge-matrix} gives
\[
  e^{\mathsf T}g=e^{\mathsf T}M_Ue-q=0,
\]
so $g\in\mathcal E_0$.
The constant-one graphon has matrix $M_1=P$, and hence $M_W=P-M_U$.
Therefore, relative to $\R e\oplus\mathcal E_0$,
\begin{equation}\label{eq:block-UW}
  M_U=
  \begin{pmatrix}
    q&g^{\mathsf T}\\
    g&A
  \end{pmatrix},
  \qquad
  M_W=
  \begin{pmatrix}
    p&-g^{\mathsf T}\\
    -g&-A
  \end{pmatrix}.
\end{equation}

\begin{lemma}[Eigenvalue bound]\label{lem:half-norm}
  Every eigenvalue of $A$ lies in $[-1/2,1/2]$.
\end{lemma}

\begin{proof}
  Throughout this proof, vector coordinates and matrix entries refer to the standard basis of $\R^N$, indexed by the partition intervals $I_1,\ldots,I_N$.
  In these coordinates,
  \[
    0\le (M_U)_{ab}=\frac{U_{ab}}{N}\le\frac1N=(M_1)_{ab},
    \qquad 1\le a,b\le N,
  \]
  so $0\le M_U\le M_1$ entrywise.
  Let $f\in\mathcal E_0$ with $\norm f=1$, and define its coordinatewise
  positive and negative parts by
  \[
    (f_+)_i=\max\{f_i,0\},
    \qquad
    (f_-)_i=\max\{-f_i,0\}.
  \]
  Both vectors have nonnegative entries, and $f=f_+-f_-$ and $|f|=f_++f_-$.
  Since $e^{\mathsf T}f=0$,
  \[
    e^{\mathsf T}f_+=e^{\mathsf T}f_-=:a,
    \qquad
    2a=e^{\mathsf T}|f|\le\norm e\norm f=1.
  \]
  Expanding $f=f_+-f_-$ gives
  \[
    f^{\mathsf T}M_Uf
    =f_+^{\mathsf T}M_Uf_+
     -f_+^{\mathsf T}M_Uf_-
     -f_-^{\mathsf T}M_Uf_+
     +f_-^{\mathsf T}M_Uf_-.
  \]
  Since $f_+$, $f_-$, and $M_U$ have nonnegative entries,
  $f_+^{\mathsf T}M_Uf_-$ and $f_-^{\mathsf T}M_Uf_+$ are nonnegative.
  Dropping these terms gives
  \[
    f^{\mathsf T}M_Uf
    \le f_+^{\mathsf T}M_Uf_++f_-^{\mathsf T}M_Uf_-
      \le2a^2\le\frac12,
  \]
  where the second inequality uses the entrywise comparison $M_U\le M_1$
  in the standard coordinates, the nonnegative entries of $f_+$ and $f_-$,
  and the identities $f_+^{\mathsf T}M_1f_+=f_-^{\mathsf T}M_1f_-=a^2$.  Similarly,
  $f_+^{\mathsf T}M_Uf_+$ and $f_-^{\mathsf T}M_Uf_-$ are nonnegative,
  so dropping these terms gives
  \[
    f^{\mathsf T}M_Uf
    \ge-f_+^{\mathsf T}M_Uf_--f_-^{\mathsf T}M_Uf_+
      \ge-2a^2\ge-\frac12.
  \]
  Here the second inequality again uses the entrywise comparison $M_U\le M_1$
  in the standard coordinates and the nonnegative entries of $f_+$ and $f_-$,
  together with $f_+^{\mathsf T}M_1f_-=f_-^{\mathsf T}M_1f_+=a^2$.
  Since $f^{\mathsf T}Af=f^{\mathsf T}M_Uf$, the Rayleigh--Ritz
  characterization of the eigenvalues of the symmetric map $A$ proves the
  claim.
\end{proof}

Let $\lambda_1,\ldots,\lambda_{N-1}$ be the eigenvalues of $A$, counted
with multiplicity, and choose corresponding orthonormal eigenvectors
$u_1,\ldots,u_{N-1}$ of $\mathcal E_0$. For each $i$, let $P_i$ be the
orthogonal projection onto the one-dimensional subspace spanned by $u_i$.
Define the finite nonnegative measure
\begin{equation*}
  \mu=\sum_{i=1}^{N-1}\norm{P_ig}^2\delta_{\lambda_i}.
\end{equation*}
The measure $\mu$ is independent of the choice of orthonormal eigenvectors
within each eigenspace.  Here $\delta_{\lambda_i}$ is the unit point mass
at $\lambda_i$.  By \Cref{lem:half-norm},
\begin{equation}\label{eq:moments}
  \supp\mu\subseteq[-1/2,1/2],
  \qquad
  g^{\mathsf T}A^jg=\int\lambda^j\dd\mu(\lambda),
  \qquad
  \mu(\R)=\norm g^2.
\end{equation}
The total mass also has the graphon expression
\begin{equation}\label{eq:degree-variance-mass}
  \mu(\R)=\norm g^2
    =\tP2U-q^2=\tP2W-p^2
    =\int_0^1\left(\int_0^1W(x,y)\dd y-p\right)^2\dd x.
\end{equation}
Indeed, $M_Ue$ records the degrees of $U$ on the partition intervals,
with normalization $1/\sqrt N$, and the degree functions of $W$ and
$U$ sum to one. Thus the quantities $g^{\mathsf T}A^jg$ are moments
of a single nonnegative measure. In the next section, we integrate
polynomials against this measure to express both the weighted
cycle--path comparison and the separate complementary-path gap.
The spectral support bound is needed later for the four-edge path
comparison, but not for the positivity of the weighted polynomials.

\section{Excursion expansions}
\label{sec:excursions}

\Cref{sec:matrix-reduction} expresses cycle and path densities through
finite matrices and represents $g^{\mathsf T}A^jg$ as moments of
$\mu$. We now derive the corresponding cycle and path excursion
expansions. Combining them gives the density-weighted expansion in
\Cref{lem:weighted-expansion}, which will prove the all-density
inequality in \Cref{sec:excursion-positivity}. We also retain the
complementary-path expansion in \Cref{lem:polynomial-reduction} for
the direct short-cycle argument in \Cref{sec:short-cycles}.

To obtain these formulas, we expand matrix powers by inserting $P+Q=I$.
For an odd integer $m\ge3$, the cycle density of $U$ satisfies
\begin{equation}\label{eq:trace-projection-expansion}
  \tC mU=\Tr(M_U^m)
  =\sum_{R_1,\ldots,R_m\in\{P,Q\}}
    \Tr(R_1M_UR_2M_U\cdots R_mM_U).
\end{equation}
We associate each summand with the sequence $R_1\cdots R_m$, called a \emph{word} in $P$ and $Q$.
This word records the projection choices; the actual matrix product contains an application of $M_U$ between consecutive projections.
For trace terms, we regard the word as cyclic, so its last and first letters are consecutive.

A \emph{block} is a maximal consecutive run of the same letter.
Thus, in a word containing both $P$ and $Q$, each $Q$-block has a $P$ immediately before and after it.
For example, the segment
\[
  P\,\underbrace{QQ}_{\text{one $Q$-block}}\,P
\]
of a word corresponds to the matrix segment $P M_U Q M_U Q M_U P$.
Reading this product from right to left, the rightmost $P$ retains the component in $\R e$.
An application of $M_U$ followed by $Q$ then retains the component in $\mathcal E_0$; the next application of $M_U$ followed by $Q$ retains a component in $\mathcal E_0$ again; the final application of $M_U$ followed by $P$ returns to $\R e$.
We call this passage from $\R e$ into $\mathcal E_0$ and back an \emph{excursion}.
Using the block forms of $P$ and $Q$ in \eqref{eq:block-PQ} and of $M_U$ in \eqref{eq:block-UW}, we obtain
\[
  P M_U Q M_U Q M_U P=(g^{\mathsf T}Ag)P.
\]
Here the outward transition contributes $g$, the step within $\mathcal E_0$ contributes $A$, and the return contributes $g^{\mathsf T}$.
More generally, a block of $\beta+1$ letters $Q$ has $\beta$ steps within $\mathcal E_0$ and contributes the scalar factor $g^{\mathsf T}A^\beta g$.
By \eqref{eq:moments}, this factor is $\int\lambda^\beta\dd\mu(\lambda)$, which connects each excursion to the spectral measure from \Cref{sec:matrix-reduction}.

We organize the summands by first fixing the number $r$ of excursions in a word.
For each fixed $r$, we sum over the possible lengths of the individual $P$-blocks and $Q$-blocks.
The resulting sums of monomials are expressed using complete homogeneous symmetric polynomials, which we now define.
For variables $y_1,\ldots,y_k$, the complete homogeneous symmetric
polynomial of degree $d$ is
\begin{equation}\label{eq:h-def}
  h_d(y_1,\ldots,y_k)
  =\sum_{\substack{\gamma_1,\ldots,\gamma_k\ge0\\
                    \gamma_1+\cdots+\gamma_k=d}}
   y_1^{\gamma_1}\cdots y_k^{\gamma_k},
  \qquad h_0=1.
\end{equation}
We write $\rep ak$ in an argument list when $a$ is repeated $k$ times.

\begin{lemma}[Cycle densities by excursions]
\label{lem:cycle-excursions}
  For every odd $m\ge3$,
  \begin{align}
    \tC mU
    ={}&q^m+\Tr(A^m)
    +\sum_{r=1}^{(m-1)/2}\frac mr
      \int\!\cdots\!\int
      h_{m-2r}(\rep qr,\lambda_1,\ldots,\lambda_r)
      \prod_{i=1}^r\dd\mu(\lambda_i),
    \label{eq:cycle-U}\\
    \tC mW
    ={}&p^m-\Tr(A^m)
    +\sum_{r=1}^{(m-1)/2}\frac mr
      \int\!\cdots\!\int
      h_{m-2r}(\rep pr,-\lambda_1,\ldots,-\lambda_r)
      \prod_{i=1}^r\dd\mu(\lambda_i).
    \label{eq:cycle-W}
  \end{align}
\end{lemma}

\begin{proof}
  Consider the summands in \eqref{eq:trace-projection-expansion}, indexed by the words $R_1\cdots R_m$.
  If $R_i=P$ for
  every $i$, the summand is $q^m$; if $R_i=Q$ for every $i$, it is
  $\Tr(A^m)$.

  Now consider a word containing both $P$ and $Q$.
  Choose a boundary where $Q$ is followed by $P$ in cyclic order, and rotate the word so that this $P$ is its first letter and the preceding $Q$ is its last letter.
  Cyclicity of the trace preserves the value of the summand under this rotation.
  Let $r$ be the number of $Q$-blocks in the resulting linear word; there are also $r$ $P$-blocks.
  Since the cut is between different letters, no cyclic block is split, so $r$ is also the number of cyclic $Q$-blocks in the original word and does not depend on the chosen boundary.
  For $1\le i\le r$, let $\alpha_i+1$ be the
  number of letters in the $i$th $P$-block, and let $\beta_i+1$ be the
  number of letters in the $Q$-block that follows it. For example, the
  following cyclic word has $r=2$ and odd length $13$:
  \[
    \underbrace{PPP}_{\alpha_1=2}\,
    \underbrace{QQQ}_{\beta_1=2}\,
    \underbrace{PPPPP}_{\alpha_2=4}\,
    \underbrace{QQ}_{\beta_2=1}
  \]
  Counting the letters gives $\sum_{i=1}^r((\alpha_i+1)+(\beta_i+1))=m$.
  To evaluate the corresponding trace, use the block forms in \eqref{eq:block-PQ} and \eqref{eq:block-UW}:
  \[
    P M_U P=qP,
    \qquad
    P M_U(Q M_U)^{\beta_i+1}P
    =(g^{\mathsf T}A^{\beta_i}g)P.
  \]
  The $P$-block has $\alpha_i$ internal steps, each contributing $q$, while the following $Q$-block gives one excursion.
  Thus, the pair of blocks, ending at the next $P$, satisfies
  \[
    (P M_U)^{\alpha_i+1}(Q M_U)^{\beta_i+1}P
    =q^{\alpha_i}(g^{\mathsf T}A^{\beta_i}g)P.
  \]
  Since the normalized product begins with $P$, appending a final $P$ does not change its trace, by $P^2=P$ and cyclicity.
  Applying the displayed identity successively, from the last pair of blocks to the first, leaves the product of their scalar factors times $P$.
  Taking the trace and using $\Tr(P)=1$, we obtain the length condition and trace formula
  \begin{equation}\label{eq:excursion-length-weight}
    \sum_{i=1}^r(\alpha_i+\beta_i)=m-2r,
    \qquad
    \Tr(R_1M_UR_2M_U\cdots R_mM_U)
    =q^{\alpha_1+\cdots+\alpha_r}
     \prod_{i=1}^rg^{\mathsf T}A^{\beta_i}g.
  \end{equation}

  To sum over all words with $r$ cyclic $Q$-blocks, we count pairs consisting
  of a word $R_1\cdots R_m$, whose positions are labelled $1,\ldots,m$, and
  a marked boundary where $Q$ is followed by $P$.
  The mark specifies where to start reading the word in cyclic order, and
  hence which pair of blocks is recorded as $(\alpha_1,\beta_1)$.
  Reading from the mark begins with $P$ and ends with $Q$, so no block is
  split. Each word has exactly $r$ eligible boundaries and therefore gives
  $r$ marked words.

  Conversely, fix an ordered tuple
  $(\alpha_1,\beta_1,\ldots,\alpha_r,\beta_r)$ of nonnegative integers
  satisfying the length condition in \eqref{eq:excursion-length-weight}.
  Choose any position $j\in\{1,\ldots,m\}$ for the first $P$ of its first
  block, and fill the positions in cyclic order with $\alpha_1+1$ copies
  of $P$, then $\beta_1+1$ copies of $Q$, and so on, wrapping from position
  $m$ back to position $1$. Mark the boundary immediately before $j$.
  The block lengths are positive and sum to $m$, so every choice of $j$
  gives a word with exactly $r$ cyclic $Q$-blocks.
  Beginning with $P$ and ending with $Q$ refers to reading from the mark;
  the original trace sum allows any letters at positions $1$ and $m$.
  Thus, each tuple gives $m$ marked words. Even if two placements produce
  the same letters because the word is periodic, their marked boundaries
  occupy different labelled positions, so the marked words are distinct.

  These constructions are inverse: a marked word determines both the
  ordered tuple and the position $j$ immediately after its mark.
  Give each marked word the weight $\Tr(R_1M_U\cdots R_mM_U)$.
  By cyclicity of the trace and \eqref{eq:excursion-length-weight}, this
  weight depends only on the tuple, not on $j$.
  Summing the weights using the two descriptions therefore gives
  \begin{equation}\label{eq:U-r-excursions}
      \Bigg(
        r\!\!\!\!
        \sum_{\substack{R_1\cdots R_m\in\{P,Q\}^m\\
                       R_1\cdots R_m\text{ has }r\text{ cyclic $Q$-blocks}}}
        \Tr(R_1M_U\cdots R_mM_U)\Bigg)
      =
      \Bigg(
        m \!\!\!\!
        \sum_{\substack{\alpha_1,\beta_1,\ldots,
                          \alpha_r,\beta_r\ge0\\
                          \sum_{i=1}^r(\alpha_i+\beta_i)=m-2r}}
        q^{\alpha_1+\cdots+\alpha_r}
        \prod_{i=1}^rg^{\mathsf T}A^{\beta_i}g\Bigg).
  \end{equation}
  Dividing \eqref{eq:U-r-excursions} by $r$ gives the factor $m/r$.
  Summing over $r$ and including the two unmixed words gives
  \[
    \tC mU=q^m+\Tr(A^m)
    +\sum_{r=1}^{(m-1)/2}\frac mr
      \sum_{\substack{\alpha_1,\beta_1,\ldots,
                        \alpha_r,\beta_r\ge0\\
                        \sum_{i=1}^r(\alpha_i+\beta_i)=m-2r}}
      q^{\alpha_1+\cdots+\alpha_r}
      \prod_{i=1}^rg^{\mathsf T}A^{\beta_i}g.
  \]
  Applying the moment identity \eqref{eq:moments} now gives
  \[
    \begin{aligned}
      \tC mU={}&q^m+\Tr(A^m)\\
      &+\sum_{r=1}^{(m-1)/2}\frac mr
      \int\!\cdots\!\int
      \sum_{\substack{\alpha_1,\beta_1,\ldots,
                        \alpha_r,\beta_r\ge0\\
                        \sum_{i=1}^r(\alpha_i+\beta_i)=m-2r}}
      q^{\alpha_1+\cdots+\alpha_r}
      \lambda_1^{\beta_1}\cdots\lambda_r^{\beta_r}
      \prod_{i=1}^r\dd\mu(\lambda_i).
    \end{aligned}
  \]
  The inner sum is
  $h_{m-2r}(\rep qr,\lambda_1,\ldots,\lambda_r)$ by
  \eqref{eq:h-def}.  This proves \eqref{eq:cycle-U}.

  For $M_W$, the block identity \eqref{eq:block-UW} replaces
  $(q,g,A)$ by $(p,-g,-A)$.  The two crossing signs cancel in each block,
  while
  \[
    (-g)^{\mathsf T}(-A)^{\beta_i}(-g)
    =(-1)^{\beta_i}g^{\mathsf T}A^{\beta_i}g.
  \]
  Hence $q$ becomes $p$, every $\lambda_i$ becomes $-\lambda_i$, and
  $\Tr(A^m)$ becomes $-\Tr(A^m)$.  This proves \eqref{eq:cycle-W}.
\end{proof}

We next derive the corresponding formula for path densities, which supplies the corrections in both \eqref{eq:intro-weighted-commonality} and \eqref{eq:intro-commonality}.
Let $P_k$ denote the path with $k$ edges. Its homomorphism density is
\begin{equation*}
  \tP kU=e^{\mathsf T}M_U^ke.
\end{equation*}

\begin{lemma}[Path densities by excursions]
\label{lem:path-excursions}
  For every $k\ge0$,
  \begin{equation}\label{eq:path-expansion}
    \tP kU=q^k+
    \sum_{r=1}^{\lfloor k/2\rfloor}
      \int\!\cdots\!\int
      h_{k-2r}(\rep q{(r+1)},\lambda_1,\ldots,\lambda_r)
      \prod_{i=1}^r\dd\mu(\lambda_i).
  \end{equation}
\end{lemma}

\begin{proof}
  For $k=0$ and $k=1$, the formula gives $\tP0U=e^{\mathsf T}e=1$ and
  $\tP1U=e^{\mathsf T}M_Ue=q$, respectively, with an empty sum over $r$.
  Now let $k\ge2$. Since $Pe=e$ and $e^{\mathsf T}P=e^{\mathsf T}$,
  we may insert $P$ at both ends of the matrix expression for $\tP kU$.
  Inserting $P+Q=I$ between its $k$ copies of $M_U$ and expanding gives
  \begin{equation}\label{eq:path-projection-expansion}
    \tP kU=
    \sum_{(R_1,\ldots,R_{k-1})\in\{P,Q\}^{k-1}}
      e^{\mathsf T}P M_U R_1M_U\cdots R_{k-1}M_UPe.
  \end{equation}
  Associate each summand in \eqref{eq:path-projection-expansion} with
  the word $PR_1\cdots R_{k-1}P$.
  This word has $k+1$ letters, separated by the $k$ copies of $M_U$.
  We read it from its fixed left endpoint to its fixed right endpoint;
  its first and last letters are separate occurrences of $P$.
  If every letter is $P$, the identity $P M_U P=qP$ gives the summand
  $q^k e^{\mathsf T}Pe=q^k$.

  Now fix the number $r\ge1$ of $Q$-blocks in a word.
  Because the word starts and ends with $P$, it has $r+1$ $P$-blocks,
  alternating with the $r$ $Q$-blocks. Let
  $\alpha_0+1,\ldots,\alpha_r+1$ be the numbers of letters in the
  successive $P$-blocks, and let $\beta_1+1,\ldots,\beta_r+1$ be the
  numbers of letters in the successive $Q$-blocks.  For example, the
  following word has $r=2$:
  \[
    \underbrace{PP}_{\alpha_0=1}\,
    \underbrace{QQ}_{\beta_1=1}\,
    \underbrace{PP}_{\alpha_1=1}\,
    \underbrace{Q}_{\beta_2=0}\,
    \underbrace{PP}_{\alpha_2=1}.
  \]
  Counting the $k+1$ letters gives
  \[
    \sum_{i=0}^r(\alpha_i+1)+\sum_{i=1}^r(\beta_i+1)=k+1,
    \qquad
    \sum_{i=0}^r\alpha_i+\sum_{i=1}^r\beta_i=k-2r.
  \]
  In particular, $r\le\lfloor k/2\rfloor$.
  To evaluate the summand in \eqref{eq:path-projection-expansion}
  associated with this word, use the same block identities as in the proof
  of \Cref{lem:cycle-excursions}: each internal step in a $P$-block
  contributes $q$, and the excursion through the $i$th $Q$-block
  contributes $g^{\mathsf T}A^{\beta_i}g$.
  Collapsing these blocks leaves the product of their scalar factors
  times $P$. Since $e^{\mathsf T}Pe=1$, the summand is
  \[
    e^{\mathsf T}P M_U R_1M_U\cdots R_{k-1}M_UPe
    =q^{\alpha_0+\cdots+\alpha_r}
      \prod_{i=1}^rg^{\mathsf T}A^{\beta_i}g.
  \]

  Reading the blocks from left to right determines the ordered tuple
  $(\alpha_0,\beta_1,\alpha_1,\ldots,\beta_r,\alpha_r)$ uniquely.
  Conversely, any tuple of nonnegative integers satisfying the length
  condition constructs exactly one word: start with $\alpha_0+1$ copies
  of $P$, then $\beta_1+1$ copies of $Q$, then $\alpha_1+1$ copies of $P$,
  and continue through the final $P$-block of length $\alpha_r+1$.
  By this one-to-one correspondence, summing the terms in
  \eqref{eq:path-projection-expansion} over the admissible block-length
  tuples gives
  \begin{equation}\label{eq:path-r-excursions}
    \begin{aligned}
      &\sum_{\substack{(R_1,\ldots,R_{k-1})\in\{P,Q\}^{k-1}\\
                       PR_1\cdots R_{k-1}P
                       \text{ has }r\text{ $Q$-blocks}}}
        e^{\mathsf T}P M_U R_1M_U\cdots R_{k-1}M_UPe\\
      &\qquad\qquad\qquad {}=
      \sum_{\substack{\alpha_0,\ldots,\alpha_r,\beta_1,\ldots,\beta_r\ge0\\
                      \sum_{i=0}^r\alpha_i+
                      \sum_{i=1}^r\beta_i=k-2r}}
      q^{\alpha_0+\cdots+\alpha_r}
      \prod_{i=1}^rg^{\mathsf T}A^{\beta_i}g.
    \end{aligned}
  \end{equation}
  Summing \eqref{eq:path-r-excursions} over $r$ and including the word
  consisting only of $P$ gives
  \[
    \tP kU=q^k+
    \sum_{r=1}^{\lfloor k/2\rfloor}
      \sum_{\substack{\alpha_0,\ldots,\alpha_r,\beta_1,\ldots,\beta_r\ge0\\
                      \sum_{i=0}^r\alpha_i+
                      \sum_{i=1}^r\beta_i=k-2r}}
      q^{\alpha_0+\cdots+\alpha_r}
      \prod_{i=1}^rg^{\mathsf T}A^{\beta_i}g.
  \]
  Applying the moment identity \eqref{eq:moments} now gives
  \[
    \begin{aligned}
      \tP kU={}&q^k+
      \sum_{r=1}^{\lfloor k/2\rfloor}
      \int\!\cdots\!\int
      \sum_{\substack{\alpha_0,\ldots,\alpha_r,\beta_1,\ldots,\beta_r\ge0\\
                      \sum_{i=0}^r\alpha_i+
                      \sum_{i=1}^r\beta_i=k-2r}}
      q^{\alpha_0+\cdots+\alpha_r}
      \lambda_1^{\beta_1}\cdots\lambda_r^{\beta_r}
      \prod_{i=1}^r\dd\mu(\lambda_i).
    \end{aligned}
  \]
  The inner sum is
  $h_{k-2r}(\rep q{(r+1)},\lambda_1,\ldots,\lambda_r)$ by
  \eqref{eq:h-def}: the $r+1$ exponents $\alpha_0,\ldots,\alpha_r$
  account for the $r+1$ copies of $q$, and the exponents
  $\beta_1,\ldots,\beta_r$ belong to $\lambda_1,\ldots,\lambda_r$.
  Their total degree is $k-2r$, as required.
  This proves \eqref{eq:path-expansion}.
\end{proof}

The path expansion for $W$ follows by replacing $(q,g,A)$ with
$(p,-g,-A)$ in \Cref{lem:path-excursions}. The crossing signs again
cancel in pairs, so
\begin{equation}\label{eq:path-expansion-W}
  \tP kW=p^k+
  \sum_{r=1}^{\lfloor k/2\rfloor}
    \int\!\cdots\!\int
    h_{k-2r}(\rep p{(r+1)},-\lambda_1,\ldots,-\lambda_r)
    \prod_{i=1}^r\dd\mu(\lambda_i).
\end{equation}
We can now form the density-weighted comparison for any real
coefficient $c$. The identity below is algebraic; the restriction on
$c$ used to prove positivity will be imposed in
\Cref{sec:excursion-positivity}.

\begin{lemma}[Density-weighted excursion expansion]
\label{lem:weighted-expansion}
  Let $m\ge3$ be odd, let $c\in\R$, and, for
  $1\le r\le(m-1)/2$, set $n=m-2r$. Define
  \begin{align}
    F^{(c)}_{m,r}(q;\lambda_1,\ldots,\lambda_r)
    ={}&\frac mr\Bigl[
      h_n(\rep pr,-\lambda_1,\ldots,-\lambda_r)
      +h_n(\rep qr,\lambda_1,\ldots,\lambda_r)
      \Bigr]\notag\\
    &-cp\,h_{n-1}(\rep p{(r+1)},-\lambda_1,\ldots,-\lambda_r)
      \notag\\
    &-cq\,h_{n-1}(\rep q{(r+1)},\lambda_1,\ldots,\lambda_r),
    \label{eq:weighted-polynomial}
  \end{align}
  where $p=1-q$. Then
  \begin{equation}\label{eq:weighted-expansion}
    \begin{aligned}
      &\tC mW+\tC mU-p^m-q^m\\
      &\quad-c\Bigl[p\bigl(\tP{m-1}W-p^{m-1}\bigr)
               +q\bigl(\tP{m-1}U-q^{m-1}\bigr)\Bigr]\\
      &\qquad=\sum_{r=1}^{(m-1)/2}
        \int\!\cdots\!\int
        F^{(c)}_{m,r}(q;\lambda_1,\ldots,\lambda_r)
        \prod_{i=1}^r\dd\mu(\lambda_i).
    \end{aligned}
  \end{equation}
\end{lemma}

\begin{proof}
  Add \eqref{eq:cycle-U} and \eqref{eq:cycle-W}; the two terms
  $\Tr(A^m)$ cancel. Subtract $cp$ times
  \eqref{eq:path-expansion-W} with $k=m-1$, and $cq$ times
  \eqref{eq:path-expansion} with $k=m-1$, after removing their
  respective constant terms. All four excursion sums have the same
  index range and product measure. At excursion number $r$, the
  cycle degree is $n=m-2r$ and the path degree is $n-1$.
  Combining their integrands gives
  \eqref{eq:weighted-polynomial}, proving the identity.
\end{proof}

For the separate short-cycle argument in \Cref{sec:short-cycles},
it is useful to record the analogous expansion for the
complementary-path correction without density weights.

\begin{lemma}[Reduction to excursion polynomials]
\label{lem:polynomial-reduction}
  Let $m\ge3$ be odd. For $1\le r\le(m-1)/2$, let
  \begin{equation}\label{eq:pge23-n}
    n=m-2r,
  \end{equation}
  which is a positive odd integer, and define the excursion polynomial,
  symmetric in $\lambda_1,\ldots,\lambda_r$, by
  \begin{align}
    P_{m,r}(q;\lambda_1,\ldots,\lambda_r)
    ={}&\frac mr\Bigl[
      h_n(\rep pr,-\lambda_1,\ldots,-\lambda_r)
      +h_n(\rep qr,\lambda_1,\ldots,\lambda_r)
    \Bigr]\notag\\
    &-h_{n-1}(\rep q{(r+1)},\lambda_1,\ldots,\lambda_r).
    \label{eq:Pmr}
  \end{align}
  The sum of integrals
  \begin{equation}\label{eq:Phi-expansion}
    \Phi_m:=\sum_{r=1}^{(m-1)/2}
      \int\!\cdots\!\int
      P_{m,r}(q;\lambda_1,\ldots,\lambda_r)
      \prod_{i=1}^r\dd\mu(\lambda_i)
  \end{equation}
  is the difference between the left- and right-hand sides of
  \eqref{eq:intro-commonality}.
\end{lemma}

\begin{proof}
  Adding the cycle expansions \eqref{eq:cycle-U} and
  \eqref{eq:cycle-W} cancels $\Tr(A^m)$ with $-\Tr(A^m)$ and gives the
  following identity, with $n=m-2r$ in each summand:
  \[
    \begin{aligned}
      \tC mW+\tC mU
      ={}p^m+q^m
      &+\sum_{r=1}^{(m-1)/2}\frac mr
      \int\!\cdots\!\int
        h_n(\rep pr,-\lambda_1,\ldots,-\lambda_r)
      \prod_{i=1}^r\dd\mu(\lambda_i)\\
      &+\sum_{r=1}^{(m-1)/2}\frac mr
      \int\!\cdots\!\int
        h_n(\rep qr,\lambda_1,\ldots,\lambda_r)
      \prod_{i=1}^r\dd\mu(\lambda_i).
    \end{aligned}
  \]
  For $k=m-1$, the path expansion \eqref{eq:path-expansion} uses the
  same range $1\le r\le(m-1)/2$, and its polynomial degree is
  $(m-1)-2r=n-1$. Thus,
  \[
    \tP{m-1}U=q^{m-1}+
    \sum_{r=1}^{(m-1)/2}
      \int\!\cdots\!\int
      h_{n-1}(\rep q{(r+1)},\lambda_1,\ldots,\lambda_r)
      \prod_{i=1}^r\dd\mu(\lambda_i).
  \]
  For each fixed number $r$ of excursions, the two cycle integrals and
  the path integral use the same product measure
  $\prod_{i=1}^r\dd\mu(\lambda_i)$.
  Subtracting the path contribution from the sum of the cycle
  contributions therefore combines their integrands into
  \[
    \frac mr\Bigl[
      h_n(\rep pr,-\lambda_1,\ldots,-\lambda_r)
      +h_n(\rep qr,\lambda_1,\ldots,\lambda_r)
    \Bigr]
    -h_{n-1}(\rep q{(r+1)},\lambda_1,\ldots,\lambda_r),
  \]
  which is $P_{m,r}(q;\lambda_1,\ldots,\lambda_r)$ by
  \eqref{eq:Pmr}.  Therefore,
  \begin{equation}\label{eq:commonality-gap}
    \tC mW+\tC mU-\tP{m-1}U
    =p^m+q^m-q^{m-1}+\Phi_m.
  \end{equation}
  This identifies $\Phi_m$ as the
  difference between the two sides of \eqref{eq:intro-commonality}.
\end{proof}

For any fixed $m$ and density, pointwise nonnegativity of all the
polynomials $P_{m,r}$ on $[-1/2,1/2]^r$ is sufficient to prove
\eqref{eq:intro-commonality}, since $\mu$ is nonnegative and supported
on that interval. \Cref{sec:short-cycles} uses this sufficient
condition for $C_3,C_5,C_7$. For arbitrary odd $m$, the proof in
\Cref{sec:excursion-positivity} instead uses the weighted polynomials
$F^{(c)}_{m,r}$ and a comparison of complementary path densities.
No pointwise sign assertion for the original $P_{m,r}$ at general
cycle lengths is needed in that argument.

\section{The cycles \texorpdfstring{$C_3,C_5,C_7$}{C3, C5, C7}}
\label{sec:short-cycles}

In this section, we consider the first three odd cycles to illustrate the
general excursion argument through explicit low-degree computations. The
identities used in \Cref{lem:polynomial-reduction} do not require
$p\ge2/3$, and for these cycles, every excursion polynomial is pointwise
nonnegative when $p\ge1/2$. Consequently, the inequalities in
\Cref{thm:main,thm:commonality} hold for $m\in\{3,5,7\}$ throughout the
larger range $p\ge1/2$.

\begin{proposition}\label{prop:short-cycles}
  Suppose that $0\le q\le1/2$ and set $p=1-q$. For all $m\in\{3,5,7\}$,
  $1\le r\le(m-1)/2$, and $\lambda_1,\ldots,\lambda_r\in\R$, we have
  \[
    P_{m,r}(q;\lambda_1,\ldots,\lambda_r)\ge0.
  \]
  Consequently, the inequalities in \Cref{thm:main,thm:commonality} hold
  for $m\in\{3,5,7\}$ whenever $p\ge1/2$.
\end{proposition}

\begin{proof}
  For $m=3$, only $r=1$ occurs, and
  \[
    \begin{aligned}
      P_{3,1}(q;\lambda)
      &=3\bigl[h_1(p,-\lambda)+h_1(q,\lambda)\bigr]
        -h_0(q,q,\lambda)\\
      &=3\bigl[(p-\lambda)+(q+\lambda)\bigr]-1=2.
    \end{aligned}
  \]

  For $m=5$ and $r=1$, we have $n=3$, and hence
  \[
    \begin{aligned}
      P_{5,1}(q;\lambda)
      ={}&5\bigl[h_3(p,-\lambda)+h_3(q,\lambda)\bigr]
          -h_2(q,q,\lambda)\\
      ={}&5\bigl[(p^3+q^3)+(q^2-p^2)\lambda
                  +(p+q)\lambda^2\bigr]
         -(3q^2+2q\lambda+\lambda^2)\\
      ={}&5\bigl[(1-3q+3q^2)+(2q-1)\lambda+\lambda^2\bigr]
         -(3q^2+2q\lambda+\lambda^2)\\
      ={}&4\lambda^2+(8q-5)\lambda+12q^2-15q+5.
    \end{aligned}
  \]
  Completing squares gives
  \[
    P_{5,1}(q;\lambda)
    =4\left(\lambda+q-\frac58\right)^2
     +8\left(q-\frac58\right)^2+\frac5{16}>0.
  \]
  For $r=2$, we have $n=1$, so
  \[
    \begin{aligned}
      P_{5,2}(q;\lambda_1,\lambda_2)
      ={}&\frac52\bigl[
        h_1(p,p,-\lambda_1,-\lambda_2)
        +h_1(q,q,\lambda_1,\lambda_2)\bigr]\\
        &-h_0(q,q,q,\lambda_1,\lambda_2)\\
      ={}&\frac52\bigl[(2p-\lambda_1-\lambda_2)
                       +(2q+\lambda_1+\lambda_2)\bigr]-1=4.
    \end{aligned}
  \]
  For $m=7$ and $r=1$, we have $n=5$.  Therefore,
  \[
    \begin{aligned}
      P_{7,1}(q;\lambda)
      ={}&7\bigl[h_5(p,-\lambda)+h_5(q,\lambda)\bigr]
          -h_4(q,q,\lambda)\\
      ={}&7\bigl[p^5+q^5+(q^4-p^4)\lambda
          +(p^3+q^3)\lambda^2+(q^2-p^2)\lambda^3
          +(p+q)\lambda^4\bigr]\\
         &-(5q^4+4q^3\lambda+3q^2\lambda^2
             +2q\lambda^3+\lambda^4)\\
      ={}&7\bigl[(1-5q+10q^2-10q^3+5q^4)
          +(4q^3-6q^2+4q-1)\lambda\\
         &\qquad +(1-3q+3q^2)\lambda^2
          +(2q-1)\lambda^3+\lambda^4\bigr]\\
         &-(5q^4+4q^3\lambda+3q^2\lambda^2
             +2q\lambda^3+\lambda^4)\\
      ={}&6\lambda^4+(12q-7)\lambda^3
          +(18q^2-21q+7)\lambda^2\\
         &+(24q^3-42q^2+28q-7)\lambda
          +30q^4-70q^3+70q^2-35q+7.
    \end{aligned}
  \]
  Define
  \[
    \begin{aligned}
      D(q)&=288q^2-336q+119,\\
      C(q)&=24q^3-42q^2+28q-7,\\
      N(q)&=5184q^6-18144q^5+28602q^4-25802q^3 +13874q^2-4165q+539.
    \end{aligned}
  \]
  Direct expansion gives
  \[
    P_{7,1}(q;\lambda)=
      6\left(\lambda^2+\left(q-\frac7{12}\right)\lambda\right)^2+\frac{D(q)}{24}
      \left(\lambda+\frac{12C(q)}{D(q)}\right)^2
      +\frac{N(q)}{D(q)}.
  \]
  With $y=1/2-q\ge0$,
  \[
    D(q)=288y^2+48y+23>0
  \]
  and
  \[
    8N(q)=41472y^6+20736y^5+21456y^4+7984y^3+2032y^2+316y+11>0.
  \]
  Hence $P_{7,1}(q;\lambda)>0$.

  For $r=2$, we have $n=3$, and
  \[
    \begin{aligned}
      P_{7,2}(q;\lambda_1,\lambda_2)
      ={}&\frac72\bigl[
        h_3(p,p,-\lambda_1,-\lambda_2)
        +h_3(q,q,\lambda_1,\lambda_2)\bigr]\\
        &-h_2(q,q,q,\lambda_1,\lambda_2)\\
      ={}&\frac72\bigl[4(p^3+q^3)
        +3(q^2-p^2)(\lambda_1+\lambda_2)
        +2(\lambda_1^2+\lambda_1\lambda_2+\lambda_2^2)\bigr]\\
        &-\bigl[6q^2+3q(\lambda_1+\lambda_2)
          +(\lambda_1^2+\lambda_1\lambda_2+\lambda_2^2)\bigr]\\
      ={}&6(\lambda_1^2+\lambda_1\lambda_2+\lambda_2^2)
        +\left(18q-\frac{21}{2}\right)(\lambda_1+\lambda_2)\\
        &+36q^2-42q+14.
    \end{aligned}
  \]
  Completing squares gives
  \[
      P_{7,2}(q;\lambda_1,\lambda_2)=
      \left(\frac32(\lambda_1-\lambda_2)^2
      +\frac92\left(\lambda_1+\lambda_2+2q-\frac76\right)^2
      +18\left(q-\frac7{12}\right)^2+\frac74\right)>0.
  \]
  Finally, for $r=3$, we have $n=1$, so
  \[
    \begin{aligned}
      P_{7,3}(q;\lambda_1,\lambda_2,\lambda_3)
      ={}&\frac73\bigl[
        h_1(p,p,p,-\lambda_1,-\lambda_2,-\lambda_3)
        +h_1(q,q,q,\lambda_1,\lambda_2,\lambda_3)\bigr]-1\\
      ={}&\frac73\bigl[(3p-\lambda_1-\lambda_2-\lambda_3)
        +(3q+\lambda_1+\lambda_2+\lambda_3)\bigr]-1
      \\
      ={}&6.
    \end{aligned}
  \]

  We have proved that
  $P_{m,r}(q;\lambda_1,\ldots,\lambda_r)>0$ for every
  $0\le q\le1/2$ and every $\lambda_1,\ldots,\lambda_r\in\R$.
  Since $\mu$ is a nonnegative measure, \eqref{eq:Phi-expansion} gives
  $\Phi_m\ge0$.
  By \Cref{lem:polynomial-reduction} and the discussion following its
  statement, this proves the strengthened inequality in
  \Cref{thm:commonality} for $m\in\{3,5,7\}$ and $p\ge1/2$.
  The comparison following \Cref{thm:commonality} then gives the
  odd-cycle bound in \Cref{thm:main} for the same cycle lengths and
  edge densities.
\end{proof}

\section{Density-weighted inequalities and the dense regime}
\label{sec:excursion-positivity}

We first prove \Cref{thm:weighted-commonality} and
\eqref{eq:intro-weighted-variance} at every edge density.
The weighted excursion integrands in
\eqref{eq:weighted-polynomial} admit a positive Dirichlet/Beta
representation. We then compare complementary even-path corrections
when $p\ge2/3$, obtaining \Cref{thm:commonality} and
\Cref{thm:main}. The proof for arbitrary odd $m$ does not require
pointwise nonnegativity of the unweighted polynomials $P_{m,r}$.

\subsection{Dirichlet and Beta averaging}
\label{sec:weighted-averaging}

We begin with the elementary averaging identity used in the proof.

\begin{lemma}[Dirichlet representation]
\label{lem:pge23-dirichlet-moment}
  Let $k\ge1$ and $j\ge0$ be integers, let
  $x_1,\ldots,x_k\in\R$, and let
  $(\Theta_1,\ldots,\Theta_k)$ be uniformly distributed on the
  simplex
  \[
    \Theta_i\ge0,\qquad \sum_{i=1}^k\Theta_i=1.
  \]
  Equivalently, this vector has the
  $\operatorname{Dirichlet}(1,\ldots,1)$ distribution; for $k=1$,
  this means $\Theta_1=1$. Then
  \begin{equation}\label{eq:pge23-dirichlet-moment}
    h_j(x_1,\ldots,x_k)
      =\binom{j+k-1}{k-1}
        \E\left(\sum_{i=1}^k\Theta_i x_i\right)^j.
  \end{equation}
\end{lemma}

\begin{proof}
  The uniform distribution on the simplex has monomial moments
  \begin{equation}\label{eq:simplex-monomial-moment}
    \E\prod_{i=1}^k\Theta_i^{\alpha_i}
      =\frac{(k-1)!\prod_{i=1}^k\alpha_i!}
             {(\alpha_1+\cdots+\alpha_k+k-1)!}
      \qquad(\alpha_i\ge0).
  \end{equation}
  For completeness, these moments follow by integrating successively
  over the simplex and using
  \[
    \int_0^1x^a(1-x)^b\dd x=\frac{a!b!}{(a+b+1)!}
    \qquad(a,b\text{ nonnegative integers}).
  \]
  The latter formula follows by integration by parts, starting with
  $a=0$. For the simplex integral, conditioning on the first
  coordinate leaves a simplex scaled by $1-\Theta_1$; induction on
  $k$ gives \eqref{eq:simplex-monomial-moment}, with the normalizing
  factor $(k-1)!$. For $k=1$ the same formula is immediate.

  Expanding by the multinomial theorem and applying
  \eqref{eq:simplex-monomial-moment} gives
  \begin{align*}
    \E\left(\sum_{i=1}^k\Theta_i x_i\right)^j
      &=\sum_{\alpha_1+\cdots+\alpha_k=j}
        \frac{j!}{\prod_i\alpha_i!}
        \frac{(k-1)!\prod_i\alpha_i!}{(j+k-1)!}
        \prod_i x_i^{\alpha_i}\\
      &=\frac{j!(k-1)!}{(j+k-1)!}\,
          h_j(x_1,\ldots,x_k),
  \end{align*}
  which proves the identity.
\end{proof}

We use the Beta distribution only with positive integer parameters.
For $s,r\ge1$, its density on $(0,1)$ is
\begin{equation}\label{eq:beta-density}
  \frac{(s+r-1)!}{(s-1)!(r-1)!}\,
    x^{s-1}(1-x)^{r-1}.
\end{equation}
The following grouped version of the Dirichlet representation keeps
the spectral variables arbitrary, so no separate diagonal reduction
is needed.

\begin{lemma}[Grouped Dirichlet representation]
\label{lem:grouped-dirichlet}
  Let $s,r\ge1$ and $d\ge0$ be integers, and let
  $a,b_1,\ldots,b_r\in\R$. Let
  $X\sim\BetaLaw(s,r)$ and let
  $(\Theta_1,\ldots,\Theta_r)$ be independent of $X$ and uniformly
  distributed on the $(r-1)$-simplex. Set
  $L=\sum_{i=1}^r\Theta_i b_i$. Then
  \begin{equation}\label{eq:grouped-dirichlet}
    h_d(a^{\times s},b_1,\ldots,b_r)
      =\binom{d+s+r-1}{s+r-1}
        \E\bigl(aX+(1-X)L\bigr)^d.
  \end{equation}
\end{lemma}

\begin{proof}
  From \eqref{eq:beta-density}, for $0\le j\le d$,
  \[
    \E\bigl[X^{d-j}(1-X)^j\bigr]
      =\frac{(s+r-1)!(s+d-j-1)!(r+j-1)!}
       {(s-1)!(r-1)!(s+r+d-1)!}.
  \]
  By \Cref{lem:pge23-dirichlet-moment},
  \[
    \E L^j=\frac{h_j(b_1,\ldots,b_r)}{\binom{j+r-1}{r-1}}.
  \]
  Expand the expectation in \eqref{eq:grouped-dirichlet} by the
  binomial theorem and use independence. After multiplication by
  $\binom{d+s+r-1}{s+r-1}$, the coefficient of
  $a^{d-j}h_j(b_1,\ldots,b_r)$ is
  \[
    \binom{d+s+r-1}{s+r-1}\binom dj
    \frac{\E[X^{d-j}(1-X)^j]}{\binom{j+r-1}{r-1}}
       =\binom{d-j+s-1}{s-1}.
  \]
  On the other hand, grouping the monomials of $h_d$ by their total
  degree $j$ in the $b_i$ gives
  \[
    h_d(a^{\times s},b_1,\ldots,b_r)
      =\sum_{j=0}^d\binom{d-j+s-1}{s-1}
         a^{d-j}h_j(b_1,\ldots,b_r).
  \]
  This proves the formula, including $d=0$ and $r=1$.
\end{proof}

\subsection{Positivity of the weighted excursion integrands}

\begin{lemma}[Positive weighted expectation]
\label{lem:weighted-positivity}
  Let $m\ge3$ be odd, let $1\le r\le(m-1)/2$, and set $n=m-2r$.
  Let $p+q=1$, and let $\lambda_1,\ldots,\lambda_r\in\R$.
  Take independent random variables
  \[
    X\sim\BetaLaw(r,r),\qquad
    (\Theta_1,\ldots,\Theta_r)
       \sim\operatorname{Dirichlet}(1,\ldots,1),
  \]
  and write
  \[
    L=\sum_{i=1}^r\Theta_i\lambda_i,\qquad
    a=pX-(1-X)L,\qquad b=qX+(1-X)L.
  \]
  For every $c\in\R$,
  \begin{equation}\label{eq:weighted-positive-expectation}
    F^{(c)}_{m,r}(q;\lambda_1,\ldots,\lambda_r)
      =\frac1r\binom{m-1}{2r-1}
       \E\bigl[(m+cr-cmX)(a^n+b^n)\bigr].
  \end{equation}
  In particular,
  \begin{equation}\label{eq:weighted-pointwise-positive}
    F^{(c)}_{m,r}(q;\lambda_1,\ldots,\lambda_r)>0
       \qquad\text{if }0\le c\le\frac{m}{m-1}.
  \end{equation}
\end{lemma}

\begin{proof}
  Separating the exponent of an additional copy of a variable gives
  the recurrence
  \begin{equation}\label{eq:h-additional-coordinate}
    a\,h_{n-1}(a^{\times(r+1)},\boldsymbol b)
      =h_n(a^{\times(r+1)},\boldsymbol b)
       -h_n(a^{\times r},\boldsymbol b).
  \end{equation}
  Indeed, the monomials on the right that do not use the additional
  coordinate cancel. Factoring $a$ from each remaining monomial
  leaves all monomials of degree $n-1$ in the enlarged list. This
  proof also covers $a=0$; no division by $a$ is used.

  For this proof only, abbreviate the two sums of homogeneous
  polynomials by
  \begin{align*}
    H={}&h_n(p^{\times r},-\lambda_1,\ldots,-\lambda_r)
          +h_n(q^{\times r},\lambda_1,\ldots,\lambda_r),\\
    \widehat H={}&h_n(p^{\times(r+1)},-\lambda_1,\ldots,-\lambda_r)
          +h_n(q^{\times(r+1)},\lambda_1,\ldots,\lambda_r).
  \end{align*}
  Applying \eqref{eq:h-additional-coordinate} to both path terms
  in \eqref{eq:weighted-polynomial} gives
  \begin{equation}\label{eq:weighted-H-recurrence}
    F^{(c)}_{m,r}=\left(\frac mr+c\right)H-c\widehat H.
  \end{equation}
  Since $n+2r=m$, \Cref{lem:grouped-dirichlet} with $s=r$ gives
  \[
    H=\binom{m-1}{2r-1}\E(a^n+b^n).
  \]
  For $\widehat H$, the same lemma uses $s=r+1$, so the law of
  the constant-coordinate weight is $\BetaLaw(r+1,r)$. By
  \eqref{eq:beta-density}, its density is $2x$ times the
  $\BetaLaw(r,r)$ density. The independent variable $L$ has the
  same law in both representations. Consequently,
  \begin{align*}
    \widehat H
      &=2\binom m{2r}\E\bigl[X(a^n+b^n)\bigr]\\
      &=\frac mr\binom{m-1}{2r-1}
          \E\bigl[X(a^n+b^n)\bigr].
  \end{align*}
  Substitution in \eqref{eq:weighted-H-recurrence} proves
  \eqref{eq:weighted-positive-expectation}. All expectations are
  finite: $X\in(0,1)$ and $L$ lies between the smallest and largest
  of the finitely many fixed spectral variables.

  To prove positivity, first note that $a+b=X>0$. Since $n$ is
  positive and odd, the function $u\mapsto u^n$ is strictly
  increasing on $\R$. Thus $a>-b$ implies
  \[
    a^n+b^n>0.
  \]
  If $c=0$, the other factor in the expectation is $m>0$.
  If $0<c\le m/(m-1)$, then
  \[
    m+cr-cmX
      =m-c(m-r)+cm(1-X)>0,
  \]
  because $m-r\le m-1$, so $m-c(m-r)\ge0$, while
  $cm(1-X)>0$ almost surely. The positive constant outside the
  expectation completes the proof. Neither this argument nor the
  identity requires a spectral support bound.
\end{proof}

\begin{proposition}[Weighted inequalities with degree variation]
\label{prop:weighted-variance}
  Let $m\ge3$ be odd, let $W$ be any graphon with $\edens W=p$
  and $q=1-p$, and let $0\le c\le m/(m-1)$. Then
  \begin{equation}\label{eq:weighted-variance-general}
    \begin{aligned}
      \tC mW+\tC m{1-W}\ge{}&p^m+q^m\\
      &+c\Bigl[p\bigl(\tP{m-1}W-p^{m-1}\bigr)
         +q\bigl(\tP{m-1}{1-W}-q^{m-1}\bigr)\Bigr]\\
      &+(m-c)\bigl(\tP2W-p^2\bigr)^{(m-1)/2}.
    \end{aligned}
  \end{equation}
\end{proposition}

\begin{proof}
  First suppose that $W$ is a step graphon and use the measure
  $\mu$ of \Cref{sec:matrix-reduction}. All integrands in
  \eqref{eq:weighted-expansion} are nonnegative by
  \Cref{lem:weighted-positivity}. For $r=(m-1)/2$, we have $n=1$,
  so \eqref{eq:weighted-polynomial} gives the constant
  \[
    F^{(c)}_{m,(m-1)/2}
      =\frac mr\left(rp-\sum_i\lambda_i
                       +rq+\sum_i\lambda_i\right)-cp-cq
      =m-c.
  \]
  Retaining only this term in \eqref{eq:weighted-expansion} gives
  the lower bound $(m-c)\mu(\R)^{(m-1)/2}$ for its left-hand side.
  By \eqref{eq:degree-variance-mass},
  $\mu(\R)=\tP2W-p^2$, proving the inequality for step graphons.

  Approximate an arbitrary graphon by the same-density step
  graphons of \Cref{lem:step-reduction}. All densities in the
  inequality converge, including $\tP2W$, and $p,q$ remain fixed.
  The map $x\mapsto x^{(m-1)/2}$ is continuous on $[0,\infty)$,
  so the inequality passes to the limit.
\end{proof}

\begin{proof}[Proof of \Cref{thm:weighted-commonality} and
\eqref{eq:intro-weighted-variance}]
  Apply \Cref{prop:weighted-variance} with $c=m/(m-1)$ and discard
  its last, nonnegative term. This proves
  \eqref{eq:intro-weighted-commonality}. With $c=1$, the baseline
  terms cancel against the subtracted path baselines, giving
  \eqref{eq:intro-weighted-variance}.

  For equality in \eqref{eq:intro-weighted-commonality}, retaining
  the last term of \eqref{eq:weighted-variance-general} gives the
  strictly positive coefficient
  \[
    m-\frac{m}{m-1}=\frac{m(m-2)}{m-1}>0.
  \]
  Equality therefore forces $\tP2W=p^2$, or equivalently that
  $W$ is $p$-regular. Conversely, if $W$ is $p$-regular, then
  $\Tker W\one=p\one$ and
  $\Tker{1-W}\one=q\one$, so both path corrections vanish.
  The two self-adjoint operators preserve $\one^\perp$ and are
  negatives of each other there. Their nonconstant contributions
  to the odd traces cancel, giving
  $\tC mW+\tC m{1-W}=p^m+q^m$. The trace sums converge absolutely
  because the operators are Hilbert--Schmidt and $m\ge3$.
  Thus every regular graphon gives equality. The endpoint densities
  $p=0$ and $p=1$ are included; then $W$ is constant almost
  everywhere.
\end{proof}

\subsection{Comparison of complementary even paths}
\label{sec:dense-paths}

This is the only part of the proof that requires $p\ge2/3$.

\begin{lemma}[Complementary even-path comparison]
\label{lem:dense-path-comparison}
  Let $W$ be a graphon with $\edens W=p\ge2/3$, and set $q=1-p$.
  For every even integer $k\ge2$,
  \begin{equation}\label{eq:dense-path-comparison}
    \tP kW-p^k\ge\tP k{1-W}-q^k.
  \end{equation}
\end{lemma}

\begin{proof}
  We first work with a step graphon and the block matrices
  \eqref{eq:block-UW}. If $q=0$, then $W=1$ almost everywhere and
  both sides of \eqref{eq:dense-path-comparison} are zero.
  Hence assume $0<q\le1/3$. Recall that
  \begin{equation}\label{eq:complementary-second-moments}
    \tP2W-p^2=\norm g^2=\tP2U-q^2.
  \end{equation}
  This proves the case $k=2$.

  For $k=4$, the first columns of the squares of the block
  matrices are
  \[
    M_U^2e=\binom{q^2+\norm g^2}{qg+Ag},\qquad
    M_W^2e=\binom{p^2+\norm g^2}{-pg+Ag}.
  \]
  Since $\tP4U=\norm{M_U^2e}^2$ and
  $\tP4W=\norm{M_W^2e}^2$, expansion gives
  \begin{align*}
    \tP4W-p^4
       &=3p^2\norm g^2+\norm g^4-2p\ip g{Ag}+\norm{Ag}^2,\\
    \tP4U-q^4
       &=3q^2\norm g^2+\norm g^4+2q\ip g{Ag}+\norm{Ag}^2.
  \end{align*}
  Subtracting and using $p+q=1$ gives
  \begin{equation}\label{eq:four-path-difference}
    (\tP4W-p^4)-(\tP4U-q^4)
       =3(p-q)\norm g^2-2\ip g{Ag}.
  \end{equation}
  By \Cref{lem:half-norm},
  $\ip g{Ag}\le\tfrac12\norm g^2$, so this is at least
  $(6p-4)\norm g^2\ge0$.

  Now let $k\ge6$ be even. We compare each path correction to
  $\norm g^2$. For any real symmetric matrix $M$, choose an
  orthonormal eigenbasis $v_i$ with eigenvalues $\xi_i$, and put
  $w_i=|\ip e{v_i}|^2$. These weights are nonnegative and sum
  to one, and
  \[
    e^{\mathsf T}M^je=\sum_i w_i\xi_i^j
       \qquad(j\ge0).
  \]
  In particular, for $M=M_W$, Jensen's inequality applied to
  $u\mapsto u^{k/2}$ on $[0,\infty)$ gives
  \[
    \tP kW\ge(\tP2W)^{k/2}
       =(p^2+\norm g^2)^{k/2}.
  \]
  The binomial theorem, with $k/2$ a positive integer, yields
  \begin{equation}\label{eq:even-path-lower}
    \tP kW-p^k\ge\frac k2p^{k-2}\norm g^2.
  \end{equation}

  For the complementary matrix $M_U$, use eigenvalues $\xi_i$
  and weights $w_i$ as above. Their first and centered second
  moments satisfy
  \[
    \sum_iw_i\xi_i=q,\qquad
    \sum_iw_i(\xi_i-q)^2=\norm g^2.
  \]
  In the standard coordinates indexed by the partition,
  \[
    \|M_U\|_{\mathrm{op}}^2
      \le\Tr(M_U^2)
      =\frac1{N^2}\sum_{a,b=1}^N U_{ab}^2
      \le\frac1{N^2}\sum_{a,b=1}^N U_{ab}=q.
  \]
  Thus $|\xi_i|\le\sqrt q$ for every $i$. For any real $z$,
  the exact polynomial identity
  \begin{equation}\label{eq:power-remainder}
    z^k-q^k-kq^{k-1}(z-q)
      =(z-q)^2\sum_{j=0}^{k-2}(j+1)q^j z^{k-2-j}
  \end{equation}
  follows by multiplying out the right side. Since $q\ge0$
  and $|\xi_i|\le\sqrt q$, each summand on the right is at most
  $(j+1)q^j(\sqrt q)^{k-2-j}$. After taking the weighted sum,
  the linear term on the left vanishes. Therefore
  \begin{equation}\label{eq:even-path-upper-sum}
    \tP kU-q^k
      \le\norm g^2\sum_{j=0}^{k-2}
          (j+1)q^j(\sqrt q)^{k-2-j}.
  \end{equation}
  This uses only upper bounds on the remainder polynomial; its
  individual terms need not be nonnegative at negative eigenvalues.

  The density hypothesis gives
  \[
    \frac qp\le\frac12,\qquad
    \frac{\sqrt q}{p}\le1,\qquad
    \frac q{p^2}\le\frac34.
  \]
  For $0\le j\le k-4$, there are at least two powers of
  $\sqrt q/p$, so
  \[
    \left(\frac qp\right)^j
      \left(\frac{\sqrt q}{p}\right)^{k-2-j}
      \le\frac34\,2^{-j}.
  \]
  For the last two terms use respectively
  $(\sqrt q/p)\le1$ and $(q/p)\le1/2$. Hence
  \begin{align}
    &\frac1{p^{k-2}}\sum_{j=0}^{k-2}
         (j+1)q^j(\sqrt q)^{k-2-j}\notag\\
    &\quad\le
      \frac34\sum_{j=0}^{k-4}\frac{j+1}{2^j}
      +\frac{k-2}{2^{k-3}}+\frac{k-1}{2^{k-2}}
       =3-2^{3-k}\le\frac k2.
    \label{eq:finite-sum-bound}
  \end{align}
  For the equality, use
  $\sum_{j=0}^{k-4}(j+1)2^{-j}
    =4-(k-1)2^{-(k-4)}$; the final inequality follows from
  $k\ge6$. Combining \eqref{eq:even-path-upper-sum} and
  \eqref{eq:finite-sum-bound} gives
  \[
    \tP kU-q^k\le\frac k2p^{k-2}\norm g^2.
  \]
  Comparison with \eqref{eq:even-path-lower} proves the result
  for every even $k\ge6$.

  Finally, the same-density approximation in
  \Cref{lem:step-reduction} passes the result to arbitrary
  graphons, including the cases $k=2,4$.
\end{proof}

The Jensen step above also proves, without a density restriction,
\[
  \tP kW\ge(\tP2W)^{k/2}\ge p^k
  \qquad(k\ge2\text{ even}).
\]
If $W$ is nonregular, then $\tP2W>p^2$, so the last inequality
is strict. The same holds for $1-W$, and a nonregular graphon
has $0<p,q<1$. This verifies the strict improvement over
$p^m+q^m$ claimed after \Cref{thm:weighted-commonality}.

\begin{proof}[Proof of \Cref{thm:commonality,thm:main}]
  Let $p\ge2/3$ and apply \Cref{lem:dense-path-comparison}
  with $k=m-1$. The identity $p+q=1$ gives
  \begin{align*}
    &p\bigl(\tP{m-1}W-p^{m-1}\bigr)
       +q\bigl(\tP{m-1}{1-W}-q^{m-1}\bigr)\\
    &\qquad\ge\tP{m-1}{1-W}-q^{m-1}.
  \end{align*}
  Using \eqref{eq:intro-weighted-variance}, we obtain the
  quantitative inequality
  \begin{equation}\label{eq:dense-variance-bound}
    \begin{aligned}
      \tC mW+\tC m{1-W}\ge{}&p^m+q^m
        +\tP{m-1}{1-W}-q^{m-1}\\
        &+(m-1)\bigl(\tP2W-p^2\bigr)^{(m-1)/2}.
    \end{aligned}
  \end{equation}
  Discarding its last term proves \Cref{thm:commonality}.
  Deleting an edge from the complementary cycle gives
  $\tC m{1-W}\le\tP{m-1}{1-W}$, and hence
  \[
    \tC mW\ge p^m+q^m-q^{m-1}=p^m-pq^{m-1}.
  \]
  This proves \Cref{thm:main}. Retaining the variance term also
  gives
  \begin{equation}\label{eq:one-colour-variance-bound}
    \tC mW\ge p^m-pq^{m-1}
      +(m-1)\bigl(\tP2W-p^2\bigr)^{(m-1)/2}.
  \end{equation}
\end{proof}

\begin{remark}[Role of the density assumption]
\label{rem:pge23-threshold}
  The weighted inequalities are valid at every edge density.
  The hypothesis $p\ge2/3$ is used only in
  \Cref{lem:dense-path-comparison}: for $k=4$ it ensures
  $6p-4\ge0$, and for $k\ge6$ it supplies the three ratio bounds
  used in \eqref{eq:finite-sum-bound}. Thus the dense one-colour
  bound follows from an all-density two-colour inequality and a
  separate comparison of paths, rather than from a sign condition
  on each of the original polynomials $P_{m,r}$.
\end{remark}

\section{Equality and qualitative stability}
\label{sec:equality-stability}

The two inequalities used in the proof impose different conditions on equality.
Equality in \Cref{thm:commonality} forces constant degrees; equality in the cycle--path comparison then forces the complement to be a union of equal cliques.
Together, these conditions identify the balanced complete multipartite graphons as the only graphons attaining the bound when $2/3\le p<1$.

For a graphon $W$ of edge density $p$, define its degree function by
\[
  \dW(x)=\int_0^1 W(x,y)\dd y.
\]
We call $W$ \emph{$p$-regular} if $\dW(x)=p$ for almost every $x$.
The gap $\Phi_m$ of \eqref{eq:Phi-expansion} is a sum of integrals against the spectral measure $\mu$ of a step graphon.
With $U=1-W$ and $q=1-p$, solving \eqref{eq:commonality-gap} for $\Phi_m$ and using $q^m-q^{m-1}=-pq^{m-1}$ turns it into
\begin{equation}\label{eq:equality-gap}
  \Phi_m(W):=\tC mW+\tC mU-\tP{m-1}U-\bigl(p^m-pq^{m-1}\bigr),
\end{equation}
a combination of homomorphism densities in which no eigenvalue appears.
We take \eqref{eq:equality-gap} as the definition of the gap for an arbitrary graphon $W$; for step graphons it agrees with $\Phi_m$.
Thus, $\Phi_m(W)$ is the difference between the two sides of \eqref{eq:intro-commonality}.
The density-weighted variance estimate and the complementary-path comparison give an explicit lower bound in terms of degree variation.

\begin{proposition}[Degree variation and equality in the two-colour bound]
\label{prop:regularity-gap}
  Let $m\ge3$ be odd, and let $W$ be a graphon of edge density $p\ge2/3$.
  Then,
  \begin{equation}\label{eq:degree-variance-gap}
    \Phi_m(W)\ge(m-1)
    \left(\int_0^1(\dW(x)-p)^2\dd x\right)^{(m-1)/2}.
  \end{equation}
  In particular, equality holds in \Cref{thm:commonality} if and only if $W$ is $p$-regular.
\end{proposition}

\begin{proof}
  By \eqref{eq:dense-variance-bound},
  \[
    \Phi_m(W)\ge(m-1)\bigl(\tP2W-p^2\bigr)^{(m-1)/2}.
  \]
  Since
  \[
    \tP2W-p^2=\int_0^1(\dW(x)-p)^2\dd x,
  \]
  this is \eqref{eq:degree-variance-gap}. Here we use the weighted
  variance bound and the comparison of the two path corrections;
  no pointwise sign assertion for the original polynomials
  $P_{m,r}$ is needed.

  If $\Phi_m(W)=0$, \eqref{eq:degree-variance-gap} forces $W$ to be $p$-regular.

  Conversely, suppose that $W$ is $p$-regular.
  The degree function of $U=1-W$ is then the constant $q$, so
  \[
    \Tker W\one=\dW=p\one,
    \qquad
    \Tker U\one=q\one .
  \]
  Both operators are self-adjoint, so each leaves the orthogonal complement
  $\one^\perp=\{f\in L^2[0,1]:\int_0^1f(x)\dd x=0\}$ invariant.
  On $\one^\perp$ the two operators are negatives of each other: the identity
  $W+U\equiv1$ gives
  \[
    \Tker Wf+\Tker Uf=\Bigl(\int_0^1f(y)\dd y\Bigr)\one=\ip f\one\one
    \qquad\text{for all }f\in L^2[0,1],
  \]
  which is the orthogonal projection onto the constants and therefore vanishes on
  $\one^\perp$.
  Write $\nu_1,\nu_2,\ldots$ for the eigenvalues of the restriction of $\Tker W$ to
  $\one^\perp$.
  The eigenvalues of $\Tker W$ are then $p$ together with the $\nu_j$, and those of
  $\Tker U$ are $q$ together with the $-\nu_j$.
  Since $m$ is odd, the terms $\nu_j^m$ and $(-\nu_j)^m$ cancel in pairs, so the
  trace formula gives
  \[
    \tC mW+\tC mU
    =\Bigl(p^m+\sum_j\nu_j^m\Bigr)+\Bigl(q^m+\sum_j(-\nu_j)^m\Bigr)
    =p^m+q^m.
  \]
  Iterating $\Tker U\one=q\one$ gives $\Tker U^{m-1}\one=q^{m-1}\one$, whence
  \[
    \tP{m-1}U=\ip\one{\Tker U^{m-1}\one}=q^{m-1}.
  \]
  Substituting both evaluations in \eqref{eq:equality-gap} and using $1-p=q$,
  \[
    \Phi_m(W)
    =\bigl(p^m+q^m\bigr)-q^{m-1}-\bigl(p^m-pq^{m-1}\bigr)
    =q^m-(1-p)q^{m-1}
    =0. \qedhere
  \]
\end{proof}

The one-colour gap decomposes as
\begin{equation}\label{eq:one-colour-gap}
  \tC mW-\bigl(p^m-pq^{m-1}\bigr)
  =\Phi_m(W)+\bigl[\tP{m-1}U-\tC mU\bigr].
\end{equation}
Both terms on the right are nonnegative.
Thus, a small one-colour gap also forces small degree variation, and equality requires the cycle--path comparison to be an equality as well.

\begin{proposition}[Equality in the odd-cycle bound]
\label{prop:main-equality}
  Let $m\ge3$ be odd, and let $W$ have edge density $2/3\le p<1$.
  Equality holds in \Cref{thm:main} if and only if $p=1-1/k$ for an integer $k\ge3$ and, up to a measure-preserving relabelling and null sets, $W=W_k$.
\end{proposition}

\begin{proof}
  The converse follows from \eqref{eq:intro-kpartite-value}.
  For the forward implication, \eqref{eq:one-colour-gap} and \Cref{prop:regularity-gap} show that $W$ is $p$-regular and
  \[
    \tC mU=\tP{m-1}U=q^{m-1},\qquad q=1-p>0.
  \]
  The graphon $U$ is $q$-regular.
  Since $U\ge0$, we may split $U=U^{1/2}\cdot U^{1/2}$ and apply the Cauchy--Schwarz
  inequality to the two factors: for $f\in L^2[0,1]$ and almost every $x$,
  \[
    \bigl|\Tker Uf(x)\bigr|^2
    =\left|\int_0^1 U(x,y)^{1/2}\cdot U(x,y)^{1/2}f(y)\dd y\right|^2
    \le\left(\int_0^1 U(x,y)\dd y\right)
      \left(\int_0^1 U(x,y)\bigl|f(y)\bigr|^2\dd y\right).
  \]
  The first factor is the degree of $U$ at $x$, which equals $q$, so
  \[
    \bigl|\Tker Uf(x)\bigr|^2\le q\int_0^1 U(x,y)\bigl|f(y)\bigr|^2\dd y .
  \]
  Integrating in $x$, exchanging the order of integration, and using $q$-regularity a
  second time, now in the variable $y$,
  \[
    \norm{\Tker Uf}_2^2
    \le q\int_0^1\bigl|f(y)\bigr|^2
      \left(\int_0^1 U(x,y)\dd x\right)\dd y
    =q^2\norm f_2^2 .
  \]
  Hence $\Tker U$ has operator norm at most $q$, and being self-adjoint, every
  eigenvalue $\lambda_i$ of $\Tker U$ lies in $[-q,q]$.
  Since $m$ is odd and $0\le U\le1$,
  \begin{equation}\label{eq:regular-complement-cycle}
    q^{m-1}=\sum_i\lambda_i^m
    \le q^{m-2}\sum_i\lambda_i^2
    =q^{m-2}\intsq U(x,y)^2\ddxy
    \le q^{m-1}.
  \end{equation}
  Equality throughout implies that $U$ takes values in $\{0,1\}$ almost everywhere and every nonzero $\lambda_i$ equals $q$.
  The operator $\Tker U$ is self-adjoint and Hilbert--Schmidt.
  Since every nonzero eigenvalue equals $q$, the spectral theorem gives
  \[
    \Tker U=qP,
    \qquad P=\text{the orthogonal projection onto }\ker(\Tker U-q).
  \]
  Since $P^2=P$, this gives $\Tker U^2=q^2P^2=q^2P=q\Tker U$.
  Comparing the kernels of these two Hilbert--Schmidt operators,
  \begin{equation}\label{eq:equal-neighbourhoods}
    \int_0^1 U(x,z)U(y,z)\dd z=qU(x,y)
    \quad\text{for almost every }(x,y).
  \end{equation}

  Write $N(x)=\{y:U(x,y)=1\}$ and $|A|$ for the Lebesgue measure of $A\subseteq[0,1]$, so that $|N(x)|=q$ for almost every $x$.
  The left side of \eqref{eq:equal-neighbourhoods} is $|N(x)\cap N(y)|$ and $U$ is $\{0,1\}$-valued, so for almost every pair $(x,y)$,
  \[
    \bigl|N(x)\cap N(y)\bigr|
    =\begin{cases}
      q&\text{if }U(x,y)=1,\\
      0&\text{if }U(x,y)=0.
    \end{cases}
  \]
  Two sets of measure $q$ meeting in measure $q$ coincide, so adjacent points have equal neighbourhoods and nonadjacent points have disjoint ones.
  The remaining identities between sets and between kernels hold up to null sets, which we no longer state.

  Fix $x$ at which this dichotomy holds for almost every $y$, and put $B=N(x)$.
  Each $y\in B$ is adjacent to $x$, so $N(y)=B$; hence $U=1$ on $B\times B$, and $U=0$ between $B$ and $[0,1]\setminus B$ because $N(y)=B$ leaves $y$ no neighbour outside $B$.
  Points outside $B$ therefore keep all $q$ of their neighbours outside $B$, so the same argument applied to $[0,1]\setminus B$ extracts a further set of measure $q$, and so on.
  A remainder of measure strictly between $0$ and $q$ is impossible, since its points have degree $q$ and no neighbours in the sets already extracted.
  These sets thus partition $[0,1]$, their number $k$ satisfies $kq=1$, and $U$ is the union of the $k$ disjoint cliques $B\times B$.
  The complement of $U$ is $W_k$ after relabelling; $p\ge2/3$ implies $k\ge3$.
\end{proof}

At $p=1$, the only graphon is $W=1$ almost everywhere, which also attains equality.
This completes the uniqueness assertion in \Cref{cor:kpartite}.

To state stability, recall the cut norm and cut distance:
\[
  \norm F_\square=\sup_{S,T\subseteq[0,1]}
    \left|\int_{S\times T}F(x,y)\ddxy\right|,
  \qquad
  \delta_\square(W,V)=\inf_\varphi\norm{W-V^\varphi}_\square.
\]
Here $S,T$ are measurable, $\varphi$ ranges over measure-preserving bijections of $[0,1]$ modulo null sets, and $V^\varphi(x,y)=V(\varphi(x),\varphi(y))$.
Graphons at zero cut distance are called \emph{equivalent}.
The space of graphons modulo this equivalence is compact in cut distance, and homomorphism densities are continuous in this metric~\cite{borgs2008convergent,lovasz2012large}.

\begin{corollary}[Qualitative stability]\label{cor:stability}
  Fix an odd integer $m\ge3$ and an integer $k\ge3$, and set $p_k=1-1/k$.
  For every $\varepsilon>0$ there exists $\delta>0$ such that every graphon $W$ with
  \[
    \bigl|\edens W-p_k\bigr|\le\delta,
    \qquad
    \tC mW\le p_k^m-p_k(1-p_k)^{m-1}+\delta
  \]
  satisfies $\delta_\square(W,W_k)<\varepsilon$.
\end{corollary}

\begin{proof}
  Suppose the conclusion fails for some $\varepsilon>0$.
  Taking $\delta=1/j$ for $j\in\mathbb N$ produces graphons $W^{(j)}$ with
  \[
    \bigl|\edens{W^{(j)}}-p_k\bigr|\le\frac1j,
    \qquad
    \tC m{W^{(j)}}\le p_k^m-p_k(1-p_k)^{m-1}+\frac1j,
    \qquad
    \delta_\square(W^{(j)},W_k)\ge\varepsilon .
  \]
  By compactness, after passing to a subsequence and relabelling, $W^{(j)}$ converges in cut distance to a graphon $W$.
  Edge and cycle densities are continuous in cut distance, so the first two conditions pass to the limit as
  \[
    \edens W=p_k,
    \qquad
    \tC mW\le p_k^m-p_k(1-p_k)^{m-1}.
  \]
  Since $k\ge3$ gives $p_k\ge2/3$, \Cref{thm:main} applies to $W$ and supplies the reverse inequality, so $\tC mW=p_k^m-p_k(1-p_k)^{m-1}$; that is, equality holds in \Cref{thm:main}.
  \Cref{prop:main-equality} then makes $W$ equivalent to $W_k$.
  Finally, $\delta_\square(W^{(j)},W_k)\to\delta_\square(W,W_k)=0$, because $\delta_\square$ obeys the triangle inequality and $\delta_\square(W^{(j)},W)\to0$.
  This contradicts $\delta_\square(W^{(j)},W_k)\ge\varepsilon$.
\end{proof}

\section{Concluding remarks and further questions}
\label{sec:conclusion}

For the one-colour Goodman bound at general odd cycle lengths, the
only nontrivial range not covered in the present paper is
$p\in(1/2,2/3)$; \Cref{prop:short-cycles} already covers this range
for $C_3$, $C_5$, and $C_7$. In contrast, the density-weighted
inequality in \Cref{thm:weighted-commonality} is valid at every edge
density. To obtain the one-colour bound, we use
\Cref{lem:dense-path-comparison} to replace its weighted path
corrections by the complementary-path correction in
\Cref{thm:commonality}, and then delete one edge from the
complementary cycle. The density assumption enters only in that
path comparison. The complementary-path inequality itself can fail
when $1/2<p<2/3$, so the all-density weighted theorem does not
justify extending this deduction to the entire middle range.
The companion paper proves the original one-colour bound on $1/2<p<2/3$ by different methods; that proof has also been verified in Lean~4 and is available at \url{ADDLINK}.

For a concrete failure of the complementary-path inequality, partition
$[0,1]$ into two sets of measure $1/2$ and let $W$ have value matrix
\[
  W=\frac1{250}\begin{pmatrix}50&193\\193&65\end{pmatrix}.
\]
Then $p=501/1000$, $q=499/1000$, and exact matrix multiplication gives
\begin{equation}\label{eq:nine-cycle-counterexample}
  \tC9W+\tC9{1-W}-p^9-q^9
     -\bigl(\tP8{1-W}-q^8\bigr)
       =-\frac{227031617768589}{5\cdot10^{21}}<0.
\end{equation}
This does not contradict \Cref{thm:weighted-commonality}: its
correction uses both path densities with different coefficients.

At each density $1-1/k$ with $k\ge3$, the present argument determines both the minimum odd-cycle density and the unique extremal graphon, with qualitative stability.
The degree-variance estimate \eqref{eq:degree-variance-gap} quantifies one part of this structure: a small gap forces degrees to be nearly constant.
The subsequent compactness argument gives no explicit rate of convergence to $W_k$ in cut distance; obtaining such a rate would strengthen the structural conclusion.

The Goodman-style perspective extends naturally beyond odd cycles. For a
nonbipartite graph $H$, let $\chi_H$ be its chromatic polynomial, let
$\chi(H)$ be its chromatic number, and let $v(H)$ be its number of vertices.
Define
\[
  G_H(p):=\chi_H\left(\frac{1}{1-p}\right)(1-p)^{v(H)}.
\]
Although this formula is written using $1/(1-p)$, it is a polynomial in
$p$: expanding the chromatic polynomial cancels every denominator. At
$p=1$, we use this polynomial's continuous extension.
We say that $H$ satisfies its \emph{Goodman-style bound} if every graphon
$W$ with $\edens W=p\ge1-1/(\chi(H)-1)$ satisfies
\[
  \dens HW\ge G_H(p).
\]
The definition is motivated by the balanced complete $k$-partite graphons $W_k$:
at $p=1-1/k$,
\[
  G_H(1-1/k)=\frac{\chi_H(k)}{k^{v(H)}}=\dens H{W_k}.
\]
Consequently, if $W_k$ minimizes the $H$-density at edge density $1-1/k$
for every integer $k\ge\chi(H)-1$, then the Goodman-style bound is tight at
all of these distinguished densities. Moreover, $G_H$ is the unique
polynomial in $p$ that takes these values for every such $k$.

This leads to a broader question: which nonbipartite graphs satisfy their
Goodman-style bounds? We record several examples and obstructions whose
proofs lie beyond the scope of this paper.
\begin{itemize}
  \item If a graph is chordal and all its maximal cliques are of the same size, then it has the Goodman-style bound.
  \item The join of a clique $K_r$ and an odd cycle $C_{2k+1}$ has the
  Goodman-style bound. In particular, taking $r=1$ gives the joins
  $K_1\vee C_{2k+1}$, often called odd wheels.
  \item If the girth of a graph is even, then the graph does not have the
  Goodman-style bound.
\end{itemize}

\medskip
\noindent\textbf{AI use disclosure.}
The original proof was discovered during two weeks of interaction with
ChatGPT~5.5 and Fable~5. We subsequently revised the argument to emphasize
its combinatorial interpretation and improve its readability. The
replacement by density-weighted averaging and complementary-path
comparison was also developed with AI assistance.

\noindent\textbf{Differences between the Lean verification and the present proof.}
The formal verification reported here concerns the one-colour
lower-bound theorem. The density-weighted inequality and the
replacement proof in \Cref{sec:excursion-positivity} are not claimed
to be included in that verification; the present argument replaces
the earlier odd-to-even Gamma-moment route by Dirichlet/Beta
averaging and a complementary even-path comparison. The equality
and stability consequences in \Cref{sec:equality-stability} are also
not included in the formal verification reported here.
\tk{TODO: Add the precise Lean theorem names, repository links, and
verification coverage for the final version.}

\noindent{\textbf{Acknowledgements.}}
\tk{TODO: Add acknowledgements.}

\bibliographystyle{abbrv}
\bibliography{references}

\appendix
\crefalias{section}{appendix}

\section{Proof of the step-graphon reduction}
\label{app:step-reduction}

\begin{proof}[Proof of \Cref{lem:step-reduction}]
  For an integer $s\ge1$, set $N=2^s$. Let $I_1,\ldots,I_N$ be the
  intervals of length $1/N$, and
  define
  \[
    W_s(x,y)=N^2\int_{I_a\times I_b}W(u,v)\dd u\dd v
    \qquad ((x,y)\in I_a\times I_b).
  \]
  Each $W_s$ is a symmetric graphon with exactly the same edge density as
  $W$:
  \[
    \int_{[0,1]^2}W_s
    =\sum_{a,b=1}^N\int_{I_a\times I_b}W=p.
  \]
  Moreover,
  \begin{equation}\label{eq:step-L1}
    \norm{W_s-W}_1\longrightarrow0
    \qquad(s\to\infty).
  \end{equation}
  To verify the convergence, let $S$ be a function constant on the
  level-$s_0$ dyadic rectangles. For $s\ge s_0$, cell averaging fixes
  $S$ and is an $L^1$-contraction, so
  \[
    \norm{W_s-W}_1
    \le \norm{W_s-S}_1+\norm{S-W}_1
    \le2\norm{W-S}_1.
  \]
  Dyadic step functions are dense in $L^1([0,1]^2)$, proving
  \eqref{eq:step-L1}.

  If $F$ is a finite graph, telescoping its edge factors gives
  \[
    \left|\dens F{W_s}-\dens FW\right|
    \le |E(F)|\norm{W_s-W}_1.
  \]
  The same bound holds for $U_s=1-W_s$ and $U=1-W$. Hence all
  homomorphism densities in the stated inequalities converge, while
  $p$ and $q$ remain unchanged. The degree variance also converges,
  since it equals $\tP2{W_s}-p^2$. Applying any of these inequalities
  to each $W_s$ and passing to the limit proves the lemma.
\end{proof}

\end{document}